\section{The source integral key invariant and its analytic constancy}
\label{sec:faithful-theta-key-source}
These source proofs follow the actual Koszul invariant, including
finite nonreduced zero schemes, the same section comparisons and
the same period frames. Their construction-specific inputs are
stated where used. Integral normalization and the connected
arithmetic-family descent are separate steps; constancy alone
does not assert the final exponent-one key unit.


\section{The integral Koszul invariant behind the key rule}
\label{sec:ftb-koszul-key}
This section translates the algebraic construction used by
\srcid{ExistsIntegralKeyInvariantRuleCharZero}.
The construction is independent of a compactification.
The input is a principally polarized abelian scheme \(A/R\) of
relative dimension \(g\geq1\), a symmetric ample principal line bundle
\(L\), a generator of \(L|_0\), a fibrewise nonzero theta section
\(\vartheta\in H^0(A,L)\), and a basis \(v_1,\ldots,v_g\) of its Lie
algebra.  This is the source datum \(\operatorname{KeyIn}(R)\).
In particular the principal polarization has trivial translation
kernel.  Changing the theta section or Lie basis does not change the
rigidification.

We use the following abelian-scheme laws with their indicated scope.
For \(m>0\), \(H^0(A,L^m)\) is free of rank \(m^g\) over a local base
and its displayed bases commute with base change; positive powers have
the higher-cohomology vanishing used in their affine Cech complexes.
For \(m\geq4\) these sections separate points and first-order tangent
vectors on a field fibre.  Multiplication by \(m\) is finite locally
free of degree \(m^{2g}\), is etale if \(m\) is invertible, and has
\(m^{2g}\) points over an algebraically closed field with \(m\ne0\).
We also use the theorem of the cube, symmetric multiplication
\([m]^*L\simeq L^{m^2}\), and faithfully flat descent along \([m]\).
These are precisely the laws used below; a scheme equipped only with
an ample line bundle would not be sufficient.

\subsection{Squared torsion of the actual augmented complex}
Let \(C=\bigoplus_{j\geq0}C_j\) be a finite free graded module over a
commutative ring, with fixed finite bases, and let \(D\) have degree
\(-1\), \(D^2=0\).
A contraction is a degree \(+1\) matrix \(h\) with \(Dh+hD=1\).
When the even and odd basis sets are identified by a bijection \(e\),
put
\[
 \tau(D,h,e)
  =\det\bigl((D+h)_{\mathrm{even},e(\mathrm{even})}\bigr)^2.
 \tag{K1}\label{eq:ftb-torsion-square}
\]
The square is part of the definition and removes the ordering sign.

\begin{srlemma}[Integral contraction and unique squared torsion]
\label{lem:ftb-squared-torsion}
Over a local ring \(R\), if the residue-field complex is exact,
a contraction exists, the even and odd ranks are equal, and
\eqref{eq:ftb-torsion-square} is a unit of \(R\).
It is independent of \(h\) and \(e\).
\end{srlemma}
\begin{proof}
Over a field choose \(G\) with \(DGD=D\), by splitting
\(\operatorname{im}D\).  Exactness gives
\(N^2=0\) for \(N=1-(DG+GD)\): \(DN=ND=0\), and the image of
\(N\), being in \(\ker D=\operatorname{im}D\), is killed by \(N\).
Thus \(DG+GD\) is invertible.
Lift \(G\) entry by entry from the residue field to \(H\) over \(R\).
The matrix \(B=DH+HD\) has unit determinant and commutes with \(D\).
Then \(h_0=HB^{-1}\) satisfies \(Dh_0+h_0D=1\).
Keep only the degree \(+1\) entries of \(h_0\).  The degree-zero
part of the last identity is still \(Dh+hD=1\), and the other
entries vanish for degree reasons.

The matrix \(D+hDh\) is an involution and exchanges even and odd
degrees: \(DhD=D\), and
\(Dh(hD)=0\).  Its two off-diagonal blocks are inverse, proving
equality of the two finite ranks.
For independence take two contractions \(h_1,h_2\) and write
\(k=h_2-h_1\).  Then
\[
 Dk+kD=0,\quad (kD)^2=0,\quad
 (D+h_2)(D+h_1)=1+kD+h_2h_1 .
\]
Order the even basis by degree.  This last matrix is block triangular;
on degree \(j\) its diagonal block is \((1+kD)|_{C_j}\).
Every such block has determinant one.  For \(j=0\), the block \(kD\)
is zero.  For successive degrees use
\(\det(1+XY)=\det(1+YX)\), with
\(X=k:C_{j-1}\to C_j\) and \(Y=D:C_j\to C_{j-1}\);
the swapped product is \(-kD|_{C_{j-1}}\).
Since that product squares to zero, the determinants of
\(1+kD\) and \(1-kD\) are inverse.  Induction gives determinant one
on every degree.

Fix the same odd ordering for both contractions.
If \(a,b\) are the determinants of their even-to-odd blocks, and
\(c\) the determinant of the odd-to-even block for \(h_1\), the
previous paragraph gives \(ac=bc=1\).  Hence \(a=b\) and \(a\)
is a unit.  A different odd ordering changes a determinant only by
the sign of a permutation.  Squaring proves all assertions.
\end{proof}

Let \(s_1,\ldots,s_g\in H^0(A,L^n)\).
For an integer \(l>gn\), choose bases \(b_j\) of
\(H^0(A,L^{l-jn})\), \(0\leq j\leq g\).
The source complex has
\[
 C_{j+1}=\bigoplus_{|I|=j}H^0(A,L^{l-jn}),
 \qquad C_0=R^d,
 \tag{K2}\label{eq:ftb-koszul-complex}
\]
with the usual signed multiplication by the \(s_i\) and the final
evaluation map at \(z_1,\ldots,z_d\).
In increasing-index exterior coordinates, deletion of \(i\) has
sign \((-1)^{|\{a\in I:a<i\}|}\).
This fixes the convention used in the determinant calculations.

The source transverse-zero hypothesis means all of the following:
the \(z_a\) are sections through pairwise distinct points in
\emph{every} fibre; their images are exactly the common zeros of
the \(s_i\); a line-bundle chart is chosen at every \(z_a\);
\(s_i(z_a)=0\); and
\[
 J_a=\det\bigl(D_{v_k}s_i(z_a)\bigr)_{i,k}
                 \in R^\times .                   \tag{K3}\label{eq:ftb-jac-unit}
\]
Values and derivatives here are the coefficients in those chosen
line-bundle charts.

\begin{srlemma}[The transverse Koszul unit]\label{lem:ftb-koszul-unit}
Let \(R\) be local Noetherian and let the displayed transverse-zero
hypothesis hold.  Then the complex
\eqref{eq:ftb-koszul-complex} is exact on every algebraically closed
field fibre.  It has a unique squared torsion \(t_l\in R^\times\),
and
\[
                     d=g!\,n^g.                   \tag{K4}\label{eq:ftb-zero-count}
\]
\end{srlemma}
\begin{proof}
On a field fibre, the Jacobian condition says that the germs of the
\(s_i\) generate the maximal ideal at each listed zero.
At other points the zero ideal is the unit ideal.
The local Koszul complex is therefore exact.
Here is the source passage to global sections: refine to a finite
affine cover trivializing \(L\), use its finite-intersection Cech
cube, and successively cut by the \(s_i\).
Weak regularity on each affine intersection gives the local
division step.  The Cech complexes of the powers
\(L^{l-jn}\) have no positive cohomology because each exponent is
positive.  Induction through these cuts shows (i) exactness at every
positive Koszul degree, (ii) interpolation of arbitrary residues at
the finite zeros, and (iii) that a section vanishing at those zeros
is \(\sum_i s_i a_i\) with \(a_i\in H^0(A,L^{l-n})\).
For (ii), first interpolate the prescribed residues separately on
affine charts; their differences on overlaps lie in the ideal
\((s_1,\ldots,s_g)\), and the preceding Cech division steps glue
them.  For (iii), the same division steps glue local expressions
in that ideal.  These three assertions give exactness of the
augmented evaluation complex, including its degree-zero term.
This is the source Cech/Koszul argument, with the exponents
\(l-jn>0\) retained.

All matrices and evaluation charts commute with field base change.
Thus exactness over the algebraic closure of the residue field
descends to the residue field (apply a linear retraction to a
solution of each finite linear system).
Lemma~\ref{lem:ftb-squared-torsion} supplies \(t_l\).
Equality of the even and odd ranks gives
\[
 d=\sum_{j=0}^g(-1)^j\binom gj(l-jn)^g=g!\,n^g.
\]
The last identity is the \(g\)-th finite difference of \(X^g\);
lower powers cancel and the leading power contributes \(g!n^g\).
\end{proof}

\subsection{The split torsion determinant}
For each positive integer \(m\) choose an integral matrix
\(\alpha_m\) of order four satisfying
\(\alpha_m^t\alpha_m=mI_4\).
It can be chosen for every \(m\): write \(m=a^2+b^2+c^2+d^2\)
and use
\[
 \alpha_m=
 \begin{pmatrix}
 a&-b&-c&-d\\ b&a&-d&c\\
 c&d&a&-b\\d&-c&b&a
 \end{pmatrix}.
 \tag{K5}\label{eq:ftb-quaternion}
\]
The relation also gives \(\alpha_m\alpha_m^t=mI_4\).

Suppose \(R\) is a local Noetherian principal ideal domain,
\(N\in R^\times\), and a fibrewise distinct full \(N\)-torsion enumeration
\(\xi_1,\ldots,\xi_{N^{2g}}\) is supplied, with \(m\mid N\).
In particular \(m\) is a unit.  In the application \(N=D\),
which is explicitly assumed invertible.
The kernel of \(\alpha_m^t:A^4\to A^4\) has
\(m^{4g}\) elements in a geometric fibre.
One sees the number by Smith normal form:
if \(d_i\) are its four invariant factors, then \(d_i\mid m\),
\(\prod_i d_i=|\det\alpha_m|=m^2\), and the kernel has
\(\prod_i d_i^{2g}=m^{4g}\) points.
Choose an enumeration
\[
 \zeta_I=(\zeta_{I1},\ldots,\zeta_{I4}),
 \qquad I\in\{1,\ldots,m^g\}^{\,4}.
\]
These are actual tuples of the supplied sections.
Indeed every such kernel tuple is \(m\)-torsion, since
\(\alpha_m\alpha_m^t=mI\), and hence is \(N\)-torsion.
Reduction maps the supplied \(N\)-torsion bijectively to the
geometric \(N\)-torsion.  A relation true after reduction is true
over \(R\): two sections of the finite etale \(N\)-torsion scheme
that agree at the closed point agree on the connected local base.
Thus
\[
 \sum_r\alpha_{m,rc}\zeta_{Ir}=0\quad(c=1,\ldots,4)
 \tag{K6}\label{eq:ftb-kernel-tuples}
\]
and distinct \(I\)'s stay distinct in every fibre.

Let \(a_{Ir}:A^4\to A\) be
\(a_{Ir}(x)=\sum_c\alpha_{m,rc}x_c+\zeta_{Ir}\).
The theorem of the cube gives a generator \(w_I\) of
\[
 \left(\bigotimes_r a_{Ir}^*L\right)
       \otimes\left(\bigotimes_c\operatorname{pr}_c^*L^{-m}\right).
 \tag{K7}\label{eq:ftb-box-generator}
\]
For precision, in the Picard group the quadratic law for a symmetric
line bundle cancels all off-diagonal terms by
\(\alpha_m^t\alpha_m=mI\).
The linear cross terms in the translations cancel by
\eqref{eq:ftb-kernel-tuples}.  Constant line bundles pulled from
\(\operatorname{Spec}R\) are free because \(R\) is local.
These are exactly the terms in the source cube calculation.

Normalize a generator
\(\mu:[m]^*L\otimes L^{-m^2}\simeq\mathcal O_A\)
at the origin using the chosen rigidification.
Since each \(\zeta_{Ir}\) is \(m\)-torsion, the product of its
values and the origin value of \(w_I\) has the same line-bundle
weight as the rigidification to exponent \(4-4m^3\).
Their quotient is therefore a unit \(\rho_I\in R^\times\), namely
\[
 w_I(0)^{m^2}\prod_r\mu(\zeta_{Ir})
              =\rho_I\,\operatorname{rig}^{\,4-4m^3}.
 \tag{K8}\label{eq:ftb-rho-normalization}
\]
This identity is an equality of framed values, so is independent of
the temporary chart used to write it.

For a basis \(b_1,\ldots,b_{m^g}\) of \(H^0(A,L^m)\), divide
\(\prod_r a_{Ir}^*\vartheta\) by \(w_I\) and expand it in the
external-product basis.  Denote the coefficient matrix by \(Q_m\):
\[
 \prod_r a_{Ir}^*\vartheta
   =w_I\sum_{I'}(Q_m)_{II'}\prod_c\operatorname{pr}_c^*b_{I'_c}.
 \tag{K9}\label{eq:ftb-box-matrix}
\]
All its coefficients are in \(R\), by the relative basis and
base-change law.  The source box-product spanning proof obtains this
basis by affine base change and then one factor at a time.

\begin{srlemma}[Unit determinant, with its actual kernel argument]
\label{lem:ftb-box-det-unit}
The determinant of \(Q_m\) is a unit.  Consequently
\[
        \lambda_m=(\det Q_m)^{2m^2}\prod_I\rho_I^2
                         \in R^\times              \tag{K10}\label{eq:ftb-lambda}
\]
is defined integrally.
\end{srlemma}
\begin{proof}
It suffices to reduce to an algebraic closure \(k\) of the residue
field.  Put \(H_m=\ker(\alpha_m:A^4(k)\to A^4(k))\), and choose
one \(I_0\).  The ratios \(w_{I_0}/w_I\) transform under translation
by \(H_m\) by characters \(\chi_I\).  The numerator theta products
in \eqref{eq:ftb-box-matrix} are invariant under \(H_m\), since
all their arguments are unchanged.  Thus those products multiplied
by \(w_{I_0}/w_I\) are eigenvectors of characters \(\chi_I\).

The characters are pairwise distinct.  If two agreed, pull their
invariant ratio back along the \(r\)-th row embedding
\(A\to A^4\), \(x\mapsto(\alpha_{m,r1}x,\ldots,\alpha_{m,r4}x)\).
The relation \(\alpha_m\alpha_m^t=mI\) cancels the other rows.
After a constant trivialization this gives an isomorphism
\[
 [m]^*\bigl(t_{\zeta_{Ir}}^*L\otimes
                     t_{\zeta_{I'r}}^*L^{-1}\bigr)\simeq\mathcal O_A
\]
invariant under \(A[m]\).
Faithfully flat descent along \([m]\) gives
\(t_{\zeta_{Ir}}^*L\simeq t_{\zeta_{I'r}}^*L\).
The principal-polarization kernel is trivial, so
\(\zeta_{Ir}=\zeta_{I'r}\) for every row, hence \(I=I'\).

Each eigenvector is nonzero.  Choose a point where \(\vartheta\)
is nonzero.  The isogeny \(\alpha_m\) is surjective (its composite
with \(\alpha_m^t\) is \([m]\)); one can therefore prescribe four
nonzero theta arguments after the translation.
Eigenvectors with distinct multiplicative characters are linearly
independent over a field: in a shortest nontrivial relation,
translate by a group element distinguishing two characters and
subtract an appropriate multiple of the relation to shorten it.
A nonzero left-kernel vector of \(Q_m\), combined with
\eqref{eq:ftb-box-matrix} and multiplied by \(w_{I_0}/w_I\), would
contradict this independence.  Thus \(\det Q_m\ne0\) over \(k\).
A coefficient over the local ring with nonzero residue is a unit.
Equation~\eqref{eq:ftb-lambda} is now a product of units.
\end{proof}

\subsection{The exact formula and cancellation of the auxiliary choices}
Choose \(n,l,l'\) with \(n\geq4\), \(gn<l<l'\), and write
\[
 D=\prod_{j=0}^g(l-jn)(l'-jn),\quad s=3g+2,\quad d=g!n^g.
\]
Supply split \(D\)-torsion and the transverse tuple above, and assume
\(D\in R^\times\).  Apply
Lemma~\ref{lem:ftb-box-det-unit} to every degree \(l-jn,l'-jn\).
Let \(t_l,t_{l'}\) be the two squared Koszul units,
\(\lambda_j,\lambda'_j\) the corresponding units
\eqref{eq:ftb-lambda}, and \(J=\prod_{a=1}^dJ_a\).
Set
\[
 \begin{split}
 T_l&=t_l^{\,4D^s}
       \prod_{j=0}^g\lambda_j^{
          (-1)^j\binom gj(D/(l-jn))^s},\\
 T_{l'}&=t_{l'}^{\,4D^s}
       \prod_{j=0}^g(\lambda'_j)^{
          (-1)^j\binom gj(D/(l'-jn))^s},\\
 \kappa&=T_l^{\,gn-2l'}T_{l'}^{\,2l-gn}
                         J^{\,8D^s(l'-l)},\\
 c&=2D^s\,d(l'-l).
 \end{split}                                      \tag{K11}\label{eq:ftb-kaform}
\]
All negative powers in this formula are powers of units.

\begin{srproposition}[Integral value and its frame law]
\label{prop:ftb-integral-koszul-invariant}
The formula \eqref{eq:ftb-kaform} gives a unit of \(R\), and
commutes with local base change preserving the displayed data.
Over a characteristic-zero field, it is independent of the bases,
enumeration and charts of the fixed transverse intersection, and of
the auxiliary generators used in \eqref{eq:ftb-box-generator}.
If \(\vartheta\) is replaced by \(a\vartheta\) and the Lie basis by
the matrix \(M\), its value satisfies
\[
              a^{8c}\kappa_{\rm new}
                         =(\det M)^{4c}\kappa_{\rm old}.
 \tag{K12}\label{eq:ftb-koszul-homog}
\]
Here the tuple \(s_i\) is fixed; independence of \emph{that} tuple
is a further step, stated separately below.
\end{srproposition}
\begin{proof}
Existence over the stated local principal ideal domain follows from
the preceding lemmas, which also give \eqref{eq:ftb-zero-count}.
Base change applies to the actual multiplication and evaluation
matrices, their contractions, the cubical generators, and determinants;
the defining identities therefore base change term by term.

We record the determinant calculations, including their signs.
If a basis in degree \(m\) changes by a matrix of determinant \(b\),
its fourth external product has determinant \(b^{4m^{3g}}\).
Thus
\[
             \lambda_{\rm new}
                   =\lambda_{\rm old}b^{-8m^{3g+2}}.
\]
Rescaling a cubical generator by a unit \(z_I\) rescales row \(I\)
of \(Q_m\) by \(z_I^{-1}\), and rescales \(\rho_I\) by
\(z_I^{m^2}\); these cancel in \eqref{eq:ftb-lambda}.
A permutation of kernel tuples contributes only a determinant sign,
removed by the square.  For arbitrary two choices of kernel enumeration, first permute the
full kernel tuples to agree.  The normalized generators \(\mu\)
agree: their quotient is a base unit, determined by its origin value.
On the proper geometrically connected variety \(A^4\), corresponding
\(w_I\)'s have a constant unit ratio \(z_I\).
The origin normalization gives
\(\rho'_I=z_I^{m^2}\rho_I\); uniqueness of the external-product
coefficients gives the inverse row factor \(z_I^{-1}\) in \(Q_m\).
These factors cancel as above.  External-product independence follows
one factor at a time from independence of the original basis, so this
comparison is valid without choosing an analytic realization.

For a change of bases \(b_j\) with determinants \(b_j\), the squared
torsion changes by
\(\prod_jb_j^{2(-1)^j\binom gj}\).
If the chosen line generator at each zero changes by \(g_a\), set
\(\gamma=\prod_ag_a\), in the convention that old evaluations are
\(g_a^l\) times new evaluations in degree \(l\).
The augmented degree-zero evaluation term then gives
\[
 t_{l,\rm new}
   =t_{l,\rm old}\gamma^{-2l}
                    \prod_jb_j^{2(-1)^j\binom gj},
 \qquad J_{\rm new}=J_{\rm old}\gamma^{-gn}.
\]
The torsion comparison follows by conjugating the whole graded
matrix, with these evaluation scalars in degree zero, and taking
the determinant of the even-to-odd block.  Permuting the zero
sections again disappears after squaring.

In \(T_l\) the basis exponents cancel term by term:
\[
 4D^s\cdot2(-1)^j\binom gj
 -(D/m)^s(-1)^j\binom gj\cdot8m^s=0.
\]
Hence \(T_{l,\rm new}=T_l\gamma^{-8lD^s}\), and the analogous
formula holds for \(l'\).
The remaining exponent of \(\gamma\) in \(\kappa\) is
\[
 -8D^s\{l(gn-2l')+l'(2l-gn)+gn(l'-l)\}=0 .
\]
This proves the claimed independence for a fixed intersection.

Multiplication of \(\vartheta\) by \(a\) multiplies every row of
\(Q_m\) by \(a^4\).  There are \(m^{4g}\) rows; thus
\(\lambda_m\) is multiplied by \(a^{8m^{4g+2}}\).
The same finite-difference identity as in
\eqref{eq:ftb-zero-count} gives
\(T_l\mapsto a^{8D^sd}T_l\), and similarly for \(l'\).
The total exponent is \(-16D^sd(l'-l)=-8c\).
Changing the Lie basis multiplies each \(J_a\) by \(\det M\);
there are \(d\) zeros, giving exponent
\(8D^sd(l'-l)=4c\).
These are exactly \eqref{eq:ftb-koszul-homog}.
\end{proof}

\begin{srlemma}[Numerical data away from a residue prime]
\label{lem:ftb-key-numerical-data}
For every \(g\geq1\) and every prime \(p\), choose
\[
 n=4p,\qquad l=gn+1,\qquad l'=gn+1+p
\]
and the matrices \eqref{eq:ftb-quaternion}.
Then \(n\geq4\), \(gn<l<l'\), \(D\geq1\), \(c\geq1\),
and \(p\nmid D\).
\end{srlemma}
\begin{proof}
For \(0\leq j\leq g\),
\[
 l-jn=4p(g-j)+1,\qquad
 l'-jn=p(4(g-j)+1)+1 .
\]
Both are positive and congruent to one modulo \(p\).
Their product \(D\) is therefore positive and prime to \(p\);
the formula for \(c\) is a product of positive integers.
\end{proof}

\paragraph{What this construction supplies and what remains.}
The proposition supplies the actual integral \(\operatorname{KaVal}\)
from split torsion and split transverse zeros, including the exact
normalizations and powers.  It does not yet assert a rule independent
of the chosen sections.  The source next proves this independence by
DVR specialization along lines, first between algebraically independent
tuples and then between arbitrary admissible tuples; it then proves
holomorphy in period coordinates, normalizes the powered rule, and
descends its exponent to one by connected families.
Those are separate producer steps.  In particular neither the unit
determinant in \eqref{eq:ftb-lambda} nor the unit
\eqref{eq:ftb-kaform} alone is the exponent-one key unit used in the
carrying section.

\section{Colliding zeros, discriminants, and independence of the key invariant}
\label{sec:ftb-key-section-independence}
Retain the actual key datum and numerical parameters of
Section~\ref{sec:ftb-koszul-key}.  The section tuple is now allowed to
have nonreduced zeros in the special fibre.  The distinction matters:
the Jacobian factors need not be units individually.
The source cancels their common discriminant in the complete formula.

\subsection{The formal complete-intersection calculation}
\begin{srlemma}[The determinant class in the socle]
\label{lem:ftb-socle-determinant}
Let \(k\) be a field, \(S=k[[X_1,\ldots,X_g]]\), and
\(F_i=\sum_j a_{ij}X_j\).  Suppose every \(X_i^N\) belongs to
\((F_1,\ldots,F_g)\), for some \(N\).
Then \(\det(a_{ij})\notin(F_1,\ldots,F_g)\).
\end{srlemma}
\begin{proof}
Write \(X_i^{N+1}=\sum_jb_{ij}F_j\), put \(B=(b_{ij})\),
and let \(I=(X_i^{N+1}:1\leq i\leq g)\).
If \(\det a\in(F)\), the adjugate identity for \(B\), applied to
\(BF=(X_i^{N+1})\), gives \(\det(Ba)\in I\).
Also
\[
 (Ba)X=\operatorname{diag}(X_i^N)X .
\]
For a regular sequence \(X_1,\ldots,X_g\), two coefficient matrices
giving the same product with \(X\) have determinants equal modulo
the ideal of their common row products.  Here is the elementary
determinant argument used in the source.  Replace the rows one at
a time.  The difference row is a syzygy on the \(X_i\).
Expand the determinant difference against that row; the cofactors
\(C_j\) satisfy \(C_jX_l-C_lX_j\) in the common-product ideal,
by the adjugate identity with one row replaced by a coordinate row.
Induction on the support of a syzygy, using
\((X_j:j\in T):X_i=(X_j:j\in T)\) for \(i\notin T\), makes
its scalar product with these cofactors lie in that ideal.
Thus
\[
 \det(Ba)-\prod_iX_i^N\in I .
\]
It would follow that \(\prod_iX_i^N\in I\).
Every coefficient of a series in \(I\) whose exponent in every
variable is at most \(N\) is zero.  The coefficient of
\(\prod_iX_i^N\) in that monomial is one, a contradiction.
The colon identity above follows directly by inspecting coefficients
of \(pX_i\); it does not use an inverse to \(N!\).
\end{proof}

\begin{srlemma}[Discriminant and formal Jacobian]
\label{lem:ftb-formal-discriminant-jacobian}
Let \(R\) be local, \(S=R[[X_1,\ldots,X_g]]\), and
\(B=S/(F_1,\ldots,F_g)\) be free of finite rank over \(R\).
Assume the constant coefficients of the \(F_i\) lie in the maximal
ideal of \(R\).  For any \(R\)-basis \(b\) of \(B\),
\[
 \operatorname{disc}(b)\ \sim\
       N_{B/R}\!\left(\det(\partial_jF_i)\bmod(F)\right),
 \tag{K13}\label{eq:ftb-disc-jac}
\]
where \(\sim\) means association by an \(R\)-unit.
\end{srlemma}
\begin{proof}
The proof constructs the diagonal element used in the source.
Let
\(\mu:S\otimes_R B\to B\) be multiplication.
Its kernel is generated by
\(X_j\otimes1-1\otimes\bar X_j\).
To see this for power series, first write
\(s=C(s(0))+\sum_jX_js_j\).
If \(J\) is the ideal generated by those differences and
\(I\) is generated by \(X_j\otimes1\), this formula gives
\(\ker\mu\subseteq J+I\ker\mu\).
The kernel is finitely generated as an \(S\)-module:
\(B\) has a finite presentation over \(S\), and
\(S\otimes_R B\) is finite over \(S\).
The \(X_j\)'s are in the Jacobson radical, also after the finite
extension \(S\to S\otimes_RB\).
Nakayama therefore gives \(\ker\mu=J\).

Choose \(a_{ij}\in S\otimes_RB\) with
\[
 F_i\otimes1=\sum_j a_{ij}
              (X_j\otimes1-1\otimes\bar X_j).
\]
Map \(\det a\) to \(B\otimes_RB\), obtaining \(\Delta\).
The adjugate identity shows that
\[
       \Delta(z\otimes1-1\otimes z)=0\quad(z\in B).
 \tag{K14}\label{eq:ftb-diagonal-annihilator}
\]
Apply the derivation \(\partial_j\) in the first tensor factor and
then multiplication.  The Leibniz rule, with all diagonal
differences zero, gives
\(\mu(a_{ij})=\overline{\partial_jF_i}\).
Consequently the multiplication image of \(\Delta\) is
\[
                  \delta=\det(\partial_jF_i)\bmod(F).
\]

Write \(\Delta=\sum_{i,j}M_{ij}b_i\otimes b_j\).
We verify that \(M\) has unit determinant.
After reduction modulo the maximal ideal, the ideal generated by
the \(\bar X_j\)'s in \(B\) is nilpotent.  Indeed these elements
are in the radical of the base maximal ideal: finiteness over the
base and the Jacobson property of \(X_j\) give this radical
containment, and finite generation gives a common power.
The reduced quotient is thus supported at the origin.
Set the second tensor factor to its residue at the origin.
The image of the equations above is a matrix expression
\(\bar F_i=\sum_j\bar a_{ij}X_j\).
A common power of each \(X_j\) is in \((\bar F)\).
Lemma~\ref{lem:ftb-socle-determinant} says that the corresponding
image of \(\Delta\) is nonzero.

Here is the resulting matrix-unit argument.  Some row
\(w=\sum_jM_{i_0j}b_j\) has nonzero residue at the origin.
It is a unit of \(B\): it is a base scalar with nonzero residue
plus an element whose power lies in the base maximal ideal.
For any \(t=(t_j)\in R^{\operatorname{rank} B}\), put
\(y=w^{-1}\sum_jt_jb_j\).
Apply the linear functional \(b_{i_0}^*\) in the first factor of
\eqref{eq:ftb-diagonal-annihilator}, after multiplication by \(y\).
It gives
\[
 \sum_{i,j}M_{ij}b_{i_0}^*(b_i y)b_j=\sum_jt_jb_j .
\]
Thus right multiplication by \(M\) on row vectors is surjective.
A square matrix over a commutative ring with this property has unit
determinant.

Finally \eqref{eq:ftb-diagonal-annihilator}, with the same coordinate
functional and then summed over the basis index, gives
\[
 \sum_{i,j}M_{ij}\operatorname{Tr}_{B/R}(b_i z)b_j
                           =\delta z\quad(z\in B).
\]
If \(H_{ij}=\operatorname{Tr}_{B/R}(b_ib_j)\), the matrix of
multiplication by \(\delta\) is \(M^tH\).
Taking determinants yields
\(N_{B/R}(\delta)=\det(M)\det(H)\).
Since \(\det(M)\) is a unit, this is
\eqref{eq:ftb-disc-jac}.
\end{proof}

\begin{srlemma}[The order inside a split generic algebra]
\label{lem:ftb-formal-order-discriminant}
Let \(O\) be a complete DVR and
\(F=(F_1,\ldots,F_g)\subset O[[X_1,\ldots,X_g]]\).
Supply finitely many \(O\)-algebra maps \(e_a:O[[X]]\to O\)
with \(e_a(F_i)=0\) and \(e_a(X_i)\) in the maximal ideal.
Assume their common kernel is \((F)\), and the evaluation map
contains \(cO^d\) for a nonzero \(c\in O\).
Then \(B=O[[X]]/(F)\) has an \(O\)-basis \(\beta_1,\ldots,\beta_d\);
its point matrix \(P=(e_a(\beta_b))\) has nonzero determinant, and
\[
                 (\det P)^2\sim
                       \prod_a e_a(\det(\partial_jF_i)).
 \tag{K15}\label{eq:ftb-point-discriminant}
\]
\end{srlemma}
\begin{proof}
Evaluation embeds \(B\) in \(O^d\) and contains \(cO^d\).
Thus it is a finite torsion-free module of rank \(d\), hence free.
After passing to the fraction field the point matrix diagonalizes
all multiplication maps.
It follows that
\(\operatorname{Tr}(z)=\sum_ae_a(z)\) and
\(N(z)=\prod_ae_a(z)\).
In particular the trace Gram matrix is \(P^tP\).
The constant coefficients of \(F_i\) are in the maximal ideal,
because their value at any \(e_a\) is zero and the \(X_i\)'s
evaluate in that ideal.  If there are no points, both sides of
\eqref{eq:ftb-point-discriminant} are one.
Otherwise apply Lemma~\ref{lem:ftb-formal-discriminant-jacobian}
and the preceding trace and norm identities.
\end{proof}

\begin{srlemma}[The finite-zero Koszul statement retains multiplicities]
\label{lem:ftb-koszul-finite-zero}
Let \(k\) be an algebraically closed field, \(A/k\) the smooth
\(g\)-dimensional principally polarized abelian variety of the key
datum, \(n\geq1\), and \(s_1,\ldots,s_g\in H^0(A,L^n)\).
Suppose only that their common zero \emph{set} is finite.
For \(l>gn\), the unaugmented section Koszul complex in degrees
\(H^0(A,L^{l-jn})\) is exact in every positive exterior degree.
For arbitrary germs \(a_y\in\mathcal O_{A,y}\) at the finitely many
common zeros, a section of \(L^l\), read in the chosen local frames,
agrees with \(a_y\) modulo the full ideal \((s_1,\ldots,s_g)\).
A section whose germs belong to that ideal is a global sum
\(\sum_i s_it_i\), with \(t_i\in H^0(A,L^{l-n})\).
The ideals in these statements are not replaced by their radicals.
\end{srlemma}
\begin{proof}
We spell out why finite nonreduced zeros meet the local regularity
input used by the source.  For a prefix of \(k\) sections let \(W_k\)
be its closed zero set.  Its dimension is at most \(g-k\).
Indeed the remaining \(g-k\) sections have affine nonzero
complements (they are positive powers of an ample line on a proper
scheme), are locally single equations, and have finite common
zeros on \(W_k\).  The source's elementary dimension-cut argument
shows that a closed subset with these properties has dimension
at most the number of sections.  Here is its prime-chain mechanism:
a positive-dimensional proper closed subset cannot lie in an affine
open, since a proper affine scheme over a field is finite; thus
one can choose a specialization outside the first complement.
Cutting a prime chain by its local single equation loses at most
one step, by the principal-ideal chain theorem.
Repeat on the closed complement with the remaining sections.
After all cuts the resulting chain lies in a finite closed set,
whose points are closed because the scheme is Jacobson.
Such a chain has length zero.  Reversing the cuts bounds the
original chain by the number of equations.  Apply this to every
prefix \(W_k\).

On an affine trivializing chart, localize at a maximal ideal.
Smoothness gives a regular local ring of dimension \(g\).
If the first \(k\) section coefficients all lie in its maximal
ideal, the quotient by them has dimension at most \(g-k\),
by the preceding closed-set dimension bound and localization.
In a Cohen--Macaulay regular local ring, \(k\) elements in the
maximal ideal with this codimension form a regular sequence:
extend them to a system of parameters, whose sequence is regular,
and take its prefix.
If a prefix coefficient is a unit, its ideal is already the unit
ideal; if the next coefficient is a unit, multiplication by it
is injective on the preceding quotient.
These cases prove exactly the weak regular sequence condition
on the affine chart, with no Jacobian or reducedness premise.

Use the same finite affine Cech cube and multiplication maps as
in Lemma~\ref{lem:ftb-koszul-unit}.
The weak regularity step says that
\(s_i a\in(s_1,\ldots,s_{i-1})\) implies
\(a\in(s_1,\ldots,s_{i-1})\) on each nonempty affine intersection.
Successive cuts of the Cech cube by these multiplication maps
therefore preserve exactness modulo the successive image sums.
All the line exponents \(l-jn\) are positive, so their positive
Cech cohomology vanishes.
The elementary cube cut is explicit: lift a boundary modulo the
new image, take its differential, use weak regularity to lift
that error to the preceding cube, subtract that lift, and use
the preceding cube's exactness; the equation \(d^2=0\) makes
the subtraction consistent.
Induction over \(i\) proves positive Koszul exactness and gluing
of local classes modulo the generated ideal.

Finally prescribed germs at finitely many zeros are represented
on small affine neighbourhoods.  On overlaps their differences
are in the same ideal after the local quotient classes are fixed;
the cut-cube gluing gives a single global class.
Applied to the zero class this is the kernel assertion.
Throughout, the quotient is by \((s_i)\) itself.
Nilpotent classes in its finite local quotient are retained by
all restrictions, lifts and subtractions.
This proves all three assertions for the full finite zero scheme.
\end{proof}

\subsection{The actual global-to-formal order}
Let \(O\) be a complete DVR and let a tuple \(s_i\) on \(A/O\)
have finite common zero set on the closed fibre.
Suppose its generic zero scheme is split, reduced, and enumerated
by sections \(z_1,\ldots,z_d\), with nonzero generic Jacobians.
The sections \(z_a\) may coincide on the closed fibre.
Their distinct closed specializations form a finite set \(Y\).
Smoothness gives, at every member \(y\in Y\), the actual formal
coordinate chart \(O[[X_1,\ldots,X_g]]\to A\), obtained from an
etale coordinate neighbourhood and its completion.
Choose its coordinate origin to be one of the sections specializing
to \(y\).  Each \(z_a\) in this block is then evaluation at
\(X_i\) in the maximal ideal.  The local equations of \(s_i\)
are \(F_{y,i}\).

For a fixed \(l>gn\), let \(V_l=H^0(A,L^l)\).
The source first proves that the images of \(V_l\) generate
\[
                 B=\prod_{y\in Y}O[[X]]/(F_{y,i})
 \tag{K16}\label{eq:ftb-local-order}
\]
as an \(O\)-module; this step precedes any choice of an \(O\)-basis
of that product.  We give the two completion arguments.
On the algebraic closure of the residue field, approximate any
formal series to order \(q\) by a germ in the coordinate
neighbourhood.  The full-ideal interpolation part of
Lemma~\ref{lem:ftb-koszul-finite-zero}
produces a single section matching all these germs modulo the
equations and \((X)^q\).
Let \(\rho_1,\ldots,\rho_N\) be the images of a basis of \(V_l\).
The descending subspaces
\[
 \{v\in k^N:\ \sum_kv_k\rho_k\in(F)+(X)^q\}
\]
stabilize.  If they stabilize at \(q\), approximation at \(q+1\)
gives \((X)^q\subset(F)+(X)^{q+1}\).
Nakayama in \(k[[X]]\) gives \((X)^q\subset(F)\).
Thus the same \(\rho_k\)'s generate the exact residue quotient,
simultaneously in all blocks.

Lift this surjectivity coefficient by coefficient.
For precision the linear map has domain the product of the \(N\)
scalar coefficients and all coefficients of \(g\) multiplier series
in every block; it sends \((v,H_{yi})\) to the coefficients of
\(\sum_kv_k\rho_{y,k}+\sum_iH_{yi}F_{yi}\).
It is surjective modulo a uniformizer.
If the residue field was enlarged, apply a residue-field linear
projection to the coefficients first, and lift their projected
values to \(O\).
Successively correct modulo \(\pi,\pi^2,\ldots\).
Products of complete coefficient modules are complete, and the
target is separated; their coefficientwise limits give an exact
solution.  Consequently the map \(V_l\to B\) is surjective.

The integral kernel of ordinary point evaluation \(V_l\to O^d\)
is exactly \(\sum_i s_i H^0(A,L^{l-n})\).
Generic exactness first puts a nonzero multiple of a kernel vector
in this image.  The image is saturated: if \(\pi v=Du\), reduce
the degree-two part of \(u\) modulo \(\pi\).
The positive-degree exactness of
Lemma~\ref{lem:ftb-koszul-finite-zero}, transported to the residue
field by a linear retraction from its algebraic closure, expresses it as
\(D\bar w\).  Lift \(w\); then \(u-Dw\) is divisible by \(\pi\)
in the relevant degree, so \(v\) is again a boundary.
Induct on the power of \(\pi\) in the first nonzero multiple.
Generic surjectivity also shows that the point-evaluation image
contains \(cO^d\) for one \(c\ne0\), by clearing denominators for
the finitely many coordinate vectors.

It follows that evaluation \(B\to O^d\) is injective and contains
\(cO^d\).  Apply
Lemma~\ref{lem:ftb-formal-order-discriminant} separately to the
blocks and combine their bases.
There are matrices \(E_l\) and \(P_l\) over \(O\) such that
\[
 \operatorname{ev}_l=P_lE_l,\qquad
 E_l:V_l\twoheadrightarrow O^d,\qquad \det P_l\ne0,\qquad
       (\det P_l)^2\sim J,                          \tag{K17}\label{eq:ftb-global-order-factor}
\]
where \(J\) is the product of the Jacobians computed in the given
Lie basis and line charts.
To compare formal and Lie Jacobians, differentiate the actual
formal coordinate chart using the square-zero algebra
\(O\oplus O^g\).
The matrix of the translated Lie basis in the formal coordinate
tangents has a right inverse, by the Lie basis property and the
formal lifting property.  It therefore has unit determinant.
The chain rule multiplies each formal Jacobian by this unit.
Changing the local line generator adds a derivative of that generator
times \(s_i(z_a)=0\), so it only contributes another unit.
This proves precisely the association in
\eqref{eq:ftb-global-order-factor}, with no assumption of distinct
closed specializations.

Replacing the augmentation \(\operatorname{ev}_l\) of the Koszul
complex by \(E_l\) gives a complex with exact residue fibre.
Indeed \(E_l\) is surjective.  A residue vector in its kernel lifts
to an integral vector; subtract integral preimages of its evaluation
divided by \(\pi\), using a fixed right inverse to \(E_l\).
Its residue is unchanged and its exact evaluation is zero.
The saturated kernel calculation above then supplies the
degree-two preimage.  In higher degrees use the same residue
Koszul exactness as before.  Thus its squared torsion \(t_l^0\)
is an \(O\)-unit.  Changing the augmentation from \(E_l\) to
\(P_lE_l\) multiplies squared torsion by \((\det P_l)^2\).
This is the source finite-closed-fibre factorization, including
the integral order, not only the generic reduced zero scheme.

\subsection{Unit cancellation for the full invariant}
\begin{srproposition}[Finite closed-fibre zeros suffice]
\label{prop:ftb-key-finite-zeros-unit}
Let \(O\) be a DVR containing \(\mathbf Q\), and let \(x\) be an
actual key datum over \(O\).  Let \(s_i\) have a finite common
zero set on the closed fibre.  Let \(F\) be any field with an
injective map \(O\to F\), carrying a base change of \(x,s_i\)
and an actual value \(\kappa\) of
\eqref{eq:ftb-kaform} on a split transverse generic intersection.
If \(a\in O\), \(e\in\mathbf Z\setminus\{0\}\), and
\(\kappa^e=a\) in \(F\), then \(a\in O^\times\).
\end{srproposition}
\begin{proof}
Only finitely many points occur in the generic zero enumeration and
the two systems of split kernel tuples.
Descend their coordinates to a finitely generated subfield of \(F\)
over \(\operatorname{Frac}O\), choose a discrete valuation ring
dominating \(O\) in that field, and complete it.
The source dominating-DVR construction applies to an essentially
finite type field extension of a nonfield DVR.
Properness of \(A/O\) extends all the descended points to this
valuation ring.  A common overfield compares its completion with
\(F\); take the compositum through a maximal ideal of the tensor
product of the two fields.
This gives a complete DVR \(O'\), a local injective map \(O\to O'\),
and an injective map \(O'\to\Omega\), with all required points
integral and the original generic data preserved.
The closed-fibre zero set remains finite by quasi-finiteness of
the reduced finite closed subscheme and base change.
Because \(O\) contains \(\mathbf Q\), all positive integers,
in particular every auxiliary \(D\), remain units.

Use \eqref{eq:ftb-global-order-factor} for \(l\) and \(l'\).
After accounting for the generic line-chart factors, put
\(\gamma=\prod_a g_a\).
There are units \(t,t',\eta,\eta'\in O'^\times\) and nonzero
generic elements \(\delta,\delta'\) such that the two generic
squared Koszul torsions and the generic Jacobian product are
\[
 \begin{aligned}
 t_l&=t\,\gamma^{-2l}\delta^2,&
 t_{l'}&=t'\,\gamma^{-2l'}(\delta')^2,\\
 J_{\rm gen}&=\eta\,\gamma^{-gn}\delta^2
            =\eta'\,\gamma^{-gn}(\delta')^2 .
 \end{aligned}                                      \tag{K18}\label{eq:ftb-colliding-factors}
\]
Here \(\delta=\det P_l\) and \(\delta'=\det P_{l'}\).
The cubical determinant units \(\lambda_j,\lambda'_j\) are integral
units by Lemma~\ref{lem:ftb-box-det-unit}; the prime-to-residue
kernel tuples are fibrewise distinct, since their multiplication
maps are etale and \(O'\) is connected.

Now \((\delta')^2=(\eta/\eta')\delta^2\).
Substitute \eqref{eq:ftb-colliding-factors} into
\eqref{eq:ftb-kaform}.
The exponent of \(\gamma\) is the zero expression already computed
in Proposition~\ref{prop:ftb-integral-koszul-invariant}.
The exponent of the remaining \(\delta^2\) is
\[
 4D^{3g+2}\{(gn-2l')+(2l-gn)\}
                    +8D^{3g+2}(l'-l)=0 .
\]
Everything left is an \(O'\)-unit.  Uniqueness for the fixed generic
intersection identifies this unit with the original \(\kappa\).
Its nonzero integer power \(a\) is therefore a unit in \(O'\).
A local map reflects units, so \(a\) is a unit in \(O\).
\end{proof}

\subsection{Constancy along lines and passage to all tuples}
\begin{srlemma}[Source line argument]\label{lem:ftb-key-line}
Let \(K\) be algebraically closed of characteristic zero.
Suppose two transverse tuples \(s_0,s_1\) on a fixed key datum have
the property that the common zero set of
\((1-t)s_0+ts_1\) is finite for every \(t\in K\).
Then their values in \eqref{eq:ftb-kaform} are equal.
\end{srlemma}
\begin{proof}
Work over the parameter line \(K[T]\).
The transverse data at \(T=a\) lift after a local DVR extension of
\(K[T]_{(T-a)}\).  Indeed take a complete DVR whose residue field
contains the finitely many coordinates of the residue zeros and
prime-to-\(D\) torsion.  The Jacobian criterion makes each zero
etale; Hensel lifting gives its unique section and unit Jacobian.
Properness ensures every common zero specializes to a listed
residue zero; the local etale uniqueness therefore lists all zeros.
The split etale \(D\)-torsion lifts as well.
Proposition~\ref{prop:ftb-integral-koszul-invariant} constructs an
integral unit value there, specializing to the given one.

Field-base-change compatibility and uniqueness descend this generic
value to \(K(T)\).  Explicitly, two embeddings of an extension field
over \(K(T)\) give the same value by uniqueness.  The element is
therefore purely inseparable over \(K(T)\): its simple intermediate
field has at most one embedding into any common field.
In characteristic zero it belongs to \(K(T)\) itself.
Call it \(c(T)\ne0\).
At every \(T=a\), Proposition~\ref{prop:ftb-key-finite-zeros-unit}
says that a nonzero integer power of \(c(T)\) lying in the local
DVR is a unit.  Write \(c=P/Q\) in relatively prime polynomials.
At a root of \(P\), \(Q\) is nonzero and this criterion applied with
exponent \(1\) is a contradiction.  At a root of \(Q\), use
exponent \(-1\).
Thus \(P,Q\) have no roots in \(K\), so they are constants.
Specialization of the lifted unit then identifies both endpoint
values with this same constant.
\end{proof}

For use in this line argument, the source proves the following
finiteness statement.  Over an algebraically closed field \(k\),
let \(s\) have finite common zeros, and replace one section by the
line towards a section whose coordinates \(T_j\) in a basis of
\(H^0(L^n)\), \(n\geq4\), are algebraically independent over \(k\).
After any algebraically closed extension, every member of that
line still has finite common zeros.
Here is the proof.  Removing one equation gives a closed set \(X\)
of dimension at most one: an ample divisor on a positive-dimensional
proper component has nonempty intersection, and the principal-ideal
dimension inequality bounds that component by one when its
intersection with the original equation is finite.
If an altered zero set were infinite, a curve component would
supply two distinct closed \(k\)-points \(y_1,y_2\) outside the
original finite zero set, whose extension fibres lie in it.
(The curve is geometrically integral after selecting an irreducible
component over algebraically closed \(k\).)
Let \(a_\nu\) be the original section values and
\(\beta_{\nu j}\) the basis-section values, using line generators
at the two points.  The line equations give
\[
 (1-t)a_\nu+t\sum_jT_j\beta_{\nu j}=0 .
\]
Take \(y_1\) outside the original zero set, so \(a_1\ne0\).
Then \(t\ne0\); eliminating \(t\) and using linear independence
of the \(T_j\)'s gives
\(a_2\beta_{1j}=a_1\beta_{2j}\) for every \(j\).
Every section vanishing at \(y_1\) must therefore vanish at \(y_2\).
This contradicts point separation by \(H^0(L^n)\).
The argument permits arbitrary \(t\) in the extension field.

\begin{srproposition}[Independence of the auxiliary sections]
\label{prop:ftb-key-tuple-independent}
Fix admissible \(g,n,l,l',\alpha\) in characteristic zero.
For each key datum over an algebraically closed field \(K\), there is
one value \(\kappa_{n,l,l',\alpha}(x)\in K^\times\) such that
every actual transverse-intersection construction
\eqref{eq:ftb-kaform}, after any algebraically closed field
extension, gives its image.  This value exists on an extension,
is functorial in the framed key datum, and is independent of the
auxiliary section tuple.
\end{srproposition}
\begin{proof}
First take two coordinate tuples jointly algebraically independent
over a field of definition for \(x\) and a basis of \(H^0(L^n)\).
Generic sections meet transversally.  To verify the latter source
step, at a proposed zero descend its values and first jets to a field
of transcendence degree at most \(g\), using the dimension of \(A\).
Evaluation is one nonzero linear equation on each coefficient row.
Eliminate one coefficient per row.  The remaining determinant of
the first jets is a nonzero polynomial: first-jet separation supplies
one solution of the evaluation equations with nonzero determinant.
Its evaluation at the remaining generic coefficients is injective.
Indeed its polynomial ring has transcendence degree at most the
number of original independent coefficients, while those
coefficients belong to the image; a nonzero kernel would lower
that degree.  Thus the generic Jacobian is nonzero at every zero.
Smoothness makes these zeros isolated; properness makes their set
finite, and they are rational over the algebraically closed field.
This is the source transcendence-degree proof of generic
transversality.

Replace rows of the first generic tuple by rows of the second, one
at a time.  For each replacement, pass to the algebraic closure of
the field generated by all the retained rows.
The incoming row remains algebraically independent over this field.
The preceding finiteness argument and
Lemma~\ref{lem:ftb-key-line} give equal values at successive tuples.
Thus jointly generic tuples give the same value.
For any two separately generic tuples, adjoin a third independent
tuple; comparison through it gives the same conclusion.

Choose the universal generic tuple in
\(\Omega=\overline{K(T_j)}\).
Every \(K\)-embedding \(\Omega\to\Omega\) takes it to another
generic tuple, so fixes its invariant.
Its invariant therefore lies in \(K\).
For clarity, a nonconstant transcendental element would extend to a
transcendence basis and could be moved by the embedding replacing
each basis variable by that variable plus one; an element algebraic
over algebraically closed \(K\) is already in \(K\).

Finally let \(S_j\) be any tuple over a field extension for which an
actual transverse value is given.
Over
\(\Omega'=\overline{F(U_j)}\), use the complete DVR
\(\Omega'[[t]]\) and coordinates
\(S_j+(U_j-S_j)t\).
They specialize to \(S_j\), and are algebraically independent over
\(K\), since evaluation at \(t=1\) gives the independent \(U_j\)'s.
Lift the transverse special-fibre data by the Hensel construction
in Lemma~\ref{lem:ftb-key-line}.
Its generic value is the universal generic value; specializing the
integral unit identifies the original value with it.
Base change and transport of the framed datum follow by applying
the already proved matrix base-change law and fixed-intersection
uniqueness.  Nonzeroness holds because the generic construction is
a product of units.
\end{proof}

\paragraph{Remaining key-rule scope.}
The argument above discharges section independence and the
finite-closed-fibre unit calculation, relative only to the specified
abelian-scheme, smooth formal-coordinate, and dominating-DVR laws.
It supplies an integral algebraic invariant with the positive
exponent \(c=\operatorname{kaExp}(g,n,l,l')\).
The source still passes from this auxiliary invariant to its
normalized analytic value, replaces the auxiliary residue prime to
get all-place units, and uses connected families to obtain the
exponent-one \(\kappa\).
Those further analytic and family statements are not consequences
of section independence alone.

\paragraph{Finite-zero exactness source precision.}
The lemma used in the global-to-formal step is the more general
\srcid{KoszulSectionsExactAlgClosedPolarisedAV.lean},
\srcid{koszul\_sections\_exact\_field}
(\srcid{f0fe7fbf99f791b00095713fd70b93d5f391aec4}):
its premise is \(\mathrm{FinZ}\), not transverse zeros.
Its regularity producer is
\srcid{RingKrullDimLocalizationQuotientLe.lean}
(\srcid{3c1e8ca388fe83a5265e805d08eed06bf2788d20}),
using \srcid{TopologicalKrullDimLeOfFiniteCommonZeros.lean}
(\srcid{fdaf8d3e54271968860ca7ff57b038273c4cc344})
and \srcid{IsRegularOfFnOfRingKrullDimQuotientAddLe.lean}
(\srcid{f6e15efe6ea2c1cd4a8d9bbdbebb9f3f3ca8960f}).
The cut-cube proof is \srcid{CubeCutRegularSequence.lean}
(\srcid{ec1b32eeaff351e107dfaf7cab2a57dbb9d870ed}).
Residue exactness transport is \srcid{KosDApplyMap.lean}
(\srcid{32213dc454266e031b4d69878743e5850f8d72c0}),
and generic exactness is \srcid{KosDExactGeneric.lean}
(\srcid{77c29ad29912bd032772f06a74bd4938d6ac4ab9}).
All these proof files have the same \srcid{ThetaLawProofs/}
prefix as the other sources in this section.

\subsection{Holomorphy of the algebraic key invariant in period coordinates}
\label{sec:ftb-key-period-holomorphy}

We keep the key datum and the admissible integers \(g,n,l,l'\),
the matrices \(\alpha_m^t\alpha_m=m1_4\), and the formula
\eqref{eq:ftb-kaform}.  In this subsection a complex uniformization
means the actual isomorphism of the polarized complex abelian variety
and its rigidified line bundle with its theta torus.  In particular it
identifies algebraic sections of every positive line power with the
entire functions having the multiplier specified below.  This is the
source's \(\mathrm{ThetaUnifAV}\) section comparison, together with
the specified \(\mathrm{HolUnif}\) and \(\mathrm{Wit}\); it is a
hypothesis on the \emph{same} line bundle and frame, rather than a
choice of an unrelated analytic metric.

Write
\[
 \mathfrak H_g=\{\tau=\tau^t:\operatorname{Im}\tau>0\},\quad
 [k,n]_\tau=k+\tau n,\quad
 A_\tau(k,n;z)=
 \exp\!\left(2\pi i\left(a\cdot k-b\cdot n
                      -\tfrac12 n^t\tau n-n\cdot z\right)\right).
 \tag{H1}\label{eq:ftb-level-aut}
\]
Here \(a_i,b_i\in\{0,\tfrac12\}\).
A function of level \(m\) is an entire \(F:\mathbb C^g\to\mathbb C\)
such that
\[
 F(z+[k,n]_\tau)=A_\tau(k,n;z)^mF(z).
 \tag{H2}\label{eq:ftb-level-fn}
\]
In the uniformizing frame \(\mathrm{fr}\), an algebraic section \(s\)
of the \(m\)-th line power is read as
\[
                       s(\exp z)=F(z)\,\mathrm{fr}(z)^m.
 \tag{H3}\label{eq:ftb-fnpw}
\]
The rigidification is the one used in the cubical generators, and
the ordered Lie basis corresponds to
\((2\pi i)^{-1}e_1,\ldots,(2\pi i)^{-1}e_g\).
The parameter domain in the source consists of all complex matrices
\(p\) with \(\operatorname{Im}((p+p^t)/2)>0\);
every formula below is pulled back by \(\tau=(p+p^t)/2\).
This gives the source's open domain without treating the symmetric
matrices as an open subset of all matrices.

\begin{srlemma}[Theta bases and holomorphic coordinates]
\label{lem:ftb-holo-level-coordinates}
For every \(m\geq1\) there are \(m^g\) level-\(m\) functions
\(B_{m,c}(\tau,z)\), jointly holomorphic in \((\tau,z)\), which form
a basis at every \(\tau\in\mathfrak H_g\).
If a jointly holomorphic function belongs at each parameter to the
span of a finite jointly holomorphic, pointwise independent family,
its unique coordinates are holomorphic.
The same assertions hold for external products of the theta bases.
\end{srlemma}
\begin{proof}
With \(c\) running through representatives of
\((\mathbb Z/m\mathbb Z)^g\), use the source's basis
\[
 B_{m,c}(\tau,z)=
 \theta\!\left(\tfrac cm+a,\,mb;\,mz,m\tau\right),\qquad
 \theta(u,v;z,\tau)=
 \sum_{h\in\mathbb Z^g}
 e^{\,\pi i(h+u)^t\tau(h+u)+2\pi i(h+u)\cdot(z+v)} .
 \tag{H4}\label{eq:ftb-level-theta-basis}
\]
On a compact subset of the period and \(z\) domain the imaginary
quadratic part is bounded below by a positive multiple of
\(\lVert h\rVert^2\), while the remaining part is bounded in
absolute value by a fixed multiple of \(1+\lVert h\rVert\).
A summable Gaussian therefore dominates the terms, locally
uniformly.  The same argument after multiplication by any fixed
polynomial in \(h\) handles derivatives.  Normal convergence gives
joint holomorphy.  Changing \(h\) by an integer vector gives
\eqref{eq:ftb-level-fn}.

Here is the source's dimension argument.  First take \(a=b=0\).
The integer periods give Fourier coefficients \(c_v\), \(v\in\mathbb
Z^g\); the remaining periods imply, for each coordinate vector,
\[
 c_{v+me_j}
       =\exp(\pi i m\tau_{jj}+2\pi i\,v^t\tau e_j)c_v .
\]
Thus the \(m^g\) coefficients with residues \(0\leq v_i<m\)
determine the function: if they vanish, the recurrence makes all
Fourier coefficients vanish, and Fourier uniqueness on each real
torus makes the function zero.  Consequently the dimension is at
most \(m^g\).  Translations \(z\mapsto z+k/m\) act on the nonzero
functions in \eqref{eq:ftb-level-theta-basis} by the distinct
characters \(k\mapsto e^{2\pi i c\cdot k/m}\).
Distinct character eigenvectors are independent (apply character
independence to a purported linear relation at each \(z\)).
There are \(m^g\) of them, so they form a basis.  Multiplication by
the nonzero exponential and translation
\[
 f(z)\longmapsto
 e^{\,2\pi i m(a^t\tau a/2+a\cdot(z+b))}
                         f(z+\tau a+b)
\]
is an invertible map from level \((0,0)\) to level \((a,b)\);
it sends the displayed basis to the displayed shifted basis.

For the coordinate assertion, let \(T_1,\ldots,T_N\) be the family,
and fix a parameter \(p_0\).
The evaluation vectors \((T_j(p_0,z))_j\) span \(\mathbb C^N\);
otherwise a nonzero linear functional would be a relation among
the functions.  Choose finite linear combinations of evaluations
\(\ell_i\) with \(\ell_i(T_j(p_0,-))=\delta_{ij}\).
Set \(C_{ij}(p)=\ell_i(T_j(p,-))\) and
\(b_i(p)=\ell_i(F(p,-))\).  They are holomorphic and \(C(p_0)=1\).
On \(\det C\ne0\), the coordinate vector is \(C^{-1}b\);
the adjugate formula proves holomorphy.  Uniqueness patches these
local expressions.  Independence of external-product functions
follows one variable at a time.  To obtain their spanning
assertion, expand in the first variable, use this coordinate
argument to make the coefficients holomorphic in the remaining
variables, and iterate.  Their transformation law in the
remaining variables follows from uniqueness of the first
expansion.  This is precisely the tensor-expansion argument in
\srcid{exists\_holomorphic\_coord\_simTheta}.
\end{proof}

\begin{srlemma}[The holomorphic box determinant]
\label{lem:ftb-holo-box}
Fix \(m\geq1\), a level-\(m\) basis \(B_c\), and an integral
\(\alpha\) with \(\alpha^t\alpha=m1_4\).
The algebraic unit \(\lambda_m\) of
\eqref{eq:ftb-lambda}, read in these bases, is a holomorphic
function of the period.  Its formula is the following.
Choose integer tuples \((k_I,n_I)\) representing the two lattice
coordinates of \(\ker(\alpha^t)\), and put
\[
 t_{I,r}=\frac1m\sum_c\alpha_{rc}[k_{I,c},n_{I,c}]_\tau,\qquad
 \Phi_I(y)=
 \prod_{r=1}^4\theta(a,b;(\alpha y)_r+t_{I,r},\tau)
             e^{\,2\pi i\sum_c n_{I,c}\cdot y_c}.
\]
Write, uniquely,
\[
 \Phi_I(y)=\sum_J q_{I,J}(\tau)\prod_cB_{J_c}(\tau,y_c).
\]
Then, in the source's row convention,
\[
 \lambda_m(\tau)=
 \frac{\det(q_{I,J}(\tau))^{2m^2}}
 {\displaystyle
  \left(\prod_I\prod_r
 A_\tau\!\left(\sum_c\alpha_{rc}k_{I,c},
              \sum_c\alpha_{rc}n_{I,c};0\right)\right)^2}.
 \tag{H5}\label{eq:ftb-holo-lambda}
\]
This equality has no undetermined root of unity or scalar.
\end{srlemma}
\begin{proof}
The elementary automorphy calculation using
\(\alpha^t\alpha=m1\) and the half-characteristic gives
\[
 \Phi_I(y_1,\ldots,y_c+[k,n]_\tau,\ldots,y_4)
             =A_\tau(k,n;y_c)^m\Phi_I(y).
\]
The functions are entire and jointly holomorphic in the period.
Lemma~\ref{lem:ftb-holo-level-coordinates} therefore gives the
unique holomorphic coefficients \(q_{I,J}\).

We check the algebraic normalization.  Divide an actual cubical
generator, in each line-bundle chart, by the product of its
uniformizing frames.  On overlaps both numerator and denominator
have the same transition function; their ratio is a globally
defined nowhere-zero entire function of \(y\).
For the generator \(w_I\), its product with
\(e^{2\pi i\sum n_{I,c}\cdot y_c}\) is periodic under the full
product period lattice.  Indeed the translated frame ratios
are exactly \(A_\tau\), and the exponent cancels the
similitude cross terms.  An entire function periodic under a full
real lattice is bounded on a compact fundamental parallelepiped
and hence constant.  Call this constant \(C_I\ne0\).
This argument is a quotient of actual local generators;
it does not extend a generic rational function across a boundary.

The same argument for the multiplication generator \(\mu\) shows
that its quotient by the multiplication frames is constant.
Its origin rigidification makes that constant \(1\).
If \(Q_{I,J}\) are the algebraic box coefficients of
\eqref{eq:ftb-box-generator}, comparison with the preceding
formula gives, after the permutation of kernel representatives,
\[
 q_{I,J}=C_IQ_{I,J},\qquad
 \rho_I
 \prod_r A_\tau\!\left(\sum_c\alpha_{rc}k_{I,c},
                     \sum_c\alpha_{rc}n_{I,c};0\right)
                                  =C_I^{m^2}.
\]
Insert these into
\(\lambda_m=(\det Q)^{2m^2}\prod_I\rho_I^2\).
All \(C_I\)'s cancel.  A row permutation is removed by the even
determinant exponent.  This proves \eqref{eq:ftb-holo-lambda}.
Its denominator is an exponential and never vanishes.
\end{proof}

\begin{srlemma}[Simple zero sections and varying contractions]
\label{lem:ftb-holo-zeros-torsion}
Let \(F_1(p,z),\ldots,F_g(p,z)\) be jointly holomorphic level-\(n\)
functions.  Suppose that at \(p_0\) their common zeros on the
period torus are precisely \(d\) distinct simple zeros.
After shrinking to an open neighbourhood \(W\) of \(p_0\),
there are holomorphic lifts \(z_1(p),\ldots,z_d(p)\) of all these
zeros, still distinct modulo the periods, and all simple.
If a finite graded complex \(D(p)\) has holomorphic entries,
\(D(p)^2=0\), and is contracted at \(p_0\), its squared Koszul
torsion is a nonzero holomorphic function on a possibly smaller
neighbourhood.
\end{srlemma}
\begin{proof}
The holomorphic implicit-function theorem supplies one branch at
each old zero because its \(z\)-Jacobian is invertible.
To show that they are all the zeros, use real period coordinates
\(z=x+\tau y\) with \((x,y)\in[0,1]^{2g}\).
The complement of small neighbourhoods of the old zeros is
compact.  At each of its points at least one \(F_i(p_0,-)\)
is nonzero.  A finite subcover and continuity keep at least one
nonzero for all nearby parameters.  Inside the small
neighbourhoods implicit-function uniqueness applies.
The period lattice is discrete in these real coordinates;
shrinking again keeps the branches distinct modulo periods.
Their Jacobians remain nonzero by continuity.  Periodicity
reduces every zero to this compact parallelepiped, proving the
claimed completeness of the enumeration.

For the contraction assertion fix a degree-\(+1\) matrix \(h_0\)
with \(D(p_0)h_0+h_0D(p_0)=1\).  Put
\[
 A(p)=D(p)h_0+h_0D(p).
\]
Near \(p_0\), \(\det A(p)\ne0\).  Since \(D^2=0\), \(D\) commutes
with \(A\), and hence with \(A^{-1}\).  Thus
\(h_0A^{-1}\) is a contraction.  Retain only its degree-\(+1\)
entries to obtain \(h(p)\).  This is still a contraction:
in the anticommutator only those entries can contribute to
the degree-zero part, while all other degree parts of \(1\)
are zero.  Its entries are holomorphic by the adjugate formula.
Choose once the even--odd enumeration used at \(p_0\).
The required holomorphic function is
\[
             t(p)=\det\bigl((D(p)+h(p))_{\rm even,odd}\bigr)^2.
 \tag{H6}\label{eq:ftb-holo-torsion}
\]
The determinant is nonzero because \(D+h\) is invertible;
equivalently the graded block elimination in the squared-torsion
lemma applies.  No assertion \(h(p)^2=0\) is required.
Uniqueness of squared torsion identifies
\eqref{eq:ftb-holo-torsion} with any other contraction computation.
\end{proof}

\begin{srproposition}[The same algebraic value is locally holomorphic]
\label{prop:ftb-key-holomorphic-period}
Suppose the algebraic key datum at \(p_0\) has a value
\(\operatorname{KaVal}\) for the fixed admissible
\(n,l,l',\alpha\), and is in the normalized theta uniformization
\eqref{eq:ftb-fnpw}.
There are an open neighbourhood \(W\) of \(p_0\) and a holomorphic
function \(V\) on \(W\) such that at each \(p\in W\), \emph{every}
algebraic key datum with that normalized uniformization has
a \(\operatorname{KaVal}\) equal to \(V(p)\).
Consequently the section-independent invariant, once defined at
every period by the algebraization of its polarized theta torus,
is holomorphic in period coordinates.
\end{srproposition}
\begin{proof}
Read the initial transverse section tuple as \(F_{i,0}\).
Expand it in the level-\(n\) basis at \(p_0\), and keep its
coefficients fixed as \(p\) varies.  This gives the required
jointly holomorphic \(F_i(p,z)\) of the same level.
The preceding lemma supplies all their simple zeros \(z_e(p)\).

For \(m_j=l-jn>0\) choose the level bases
\(B_{m_j,c}(p,z)\), and likewise for \(l'\).
The products \(F_iB_{m_{j+1},c}\) have level \(m_j\).
Their multiplication coordinates are holomorphic by
Lemma~\ref{lem:ftb-holo-level-coordinates}.
The evaluations \(B_{l,c}(p,z_e(p))\) are holomorphic.
Thus all entries of the augmented Koszul matrix \(D_l(p)\)
are holomorphic, and the same holds in degree \(l'\).
Their differentials square to zero because multiplication
commutes and the evaluation points are common zeros.
At \(p_0\) the algebraic exactness already proved supplies
contractions; apply \eqref{eq:ftb-holo-torsion} in both degrees.

Let \(J(p)\) be the product of the source's normalized Jacobians
of \(F_i(p,-)\) at \(z_e(p)\); its entries are the first
\(z\)-derivatives, with the fixed \(2\pi i\) Lie normalization.
Joint holomorphy implies holomorphy of these derivatives and
their evaluations, so \(J(p)\) is holomorphic and nonzero.
Lemma~\ref{lem:ftb-holo-box} supplies the \(\lambda_j(p)\) and
\(\lambda'_j(p)\).  At \(p_0\) they are nonzero by the algebraic
box determinant theorem; shrink \(W\) so that they remain so.
Use \(t_l(p),t_{l'}(p),J(p),\lambda_j(p),\lambda'_j(p)\) in
\eqref{eq:ftb-kaform}.  Since every factor is holomorphic and
nonzero, its integer powers give a holomorphic function \(V\).

It remains to verify that this analytic formula computes the
\emph{actual} algebraic value, with no surviving chart factor.
The line-power section comparison makes the theta bases
algebraic bases and makes the \(F_i\) actual algebraic sections.
If \(U=\prod_e\mathrm{fr}_{j_e}(z_e)\), the algebraic squared
torsion in degree \(l\) is the analytic squared torsion times
\(U^{2l}\): this is conjugation of its augmented evaluation
rows by \(\operatorname{diag}(\mathrm{fr}_{j_e}(z_e)^l)\).
The corresponding algebraic Jacobian is \(J(p)U^{gn}\).
The same applies in degree \(l'\).
The total exponent of \(U\) in \eqref{eq:ftb-kaform} is
\[
 8D^{3g+2}
 \{l(gn-2l')+l'(2l-gn)+gn(l'-l)\}=0.
\]
The box factors agree by \eqref{eq:ftb-holo-lambda}, and the
squared torsions agree by their uniqueness.  Therefore the
algebraic value is exactly \(V(p)\).
Section independence proved in
Proposition~\ref{prop:ftb-key-tuple-independent} removes the
particular tuple chosen to construct \(V\).
On overlapping neighbourhoods, algebraization of a theta torus
provides a test key datum at each period; both local formulae
compute its unique algebraic value, so they agree.  This patches
the holomorphic value function.

The last algebraization clause is explicit: it is the actual
polarized-theta-torus realization, denoted
\(\mathrm{TorusAlgebraic}\) in the source.  Neither this clause
nor the section comparison is deduced from holomorphy of
\(V\).  Their match with the moduli/uniformization model is
a separate geometric constructor.
\end{proof}

\paragraph{Frozen source concordance.}
The base prefix for the following paths is
\srcid{src/Development/AbelianVarietyTheory/Proofs/}.
The analytic level-space proof is
\srcid{CovolumeBounds/EqZeroOfSemicharacterMultiplier.lean}
(\srcid{d9177a8a8151f010857a1872b029ca293d2ecdd4});
jointly holomorphic bases are
\srcid{ThetaLawProofs/ExistsLevBasisReal.lean}
(\srcid{0acd91467b163727cd9ab45b449a645d19079f73}).
The fixed-evaluation coordinate proof is
\srcid{DifferentiableOnCoordOfLinearIndependent.lean}
(\srcid{bb29ce974ad27e88014fe30ce737b1ff313a9fd8}).
The tensor expansion and its holomorphic coordinates are
\srcid{ExistsHolomorphicCoordSimilitudeThetaProduct.lean}
(\srcid{222a38b04c4c65633ddc63956623b5ebc871d2ae}).
The actual frame quotient gluing is
\srcid{ExistsEntireGeneratorDivFrames.lean}
(\srcid{5115a6c11545e30287524197a677172c8334ee56});
the full-lattice periodicity calculations are
\srcid{SimilitudeGeneratorEqExpInFrames.lean}
(\srcid{767a1a7851d799760d180c9278cc4c1670a5dfe7})
and \srcid{ThetaGeneratorPullbackMulConstInFrames.lean}
(\srcid{11d7336e26d26d4ef9f784da8f389ceab34a08c2}).
The frame automorphy identity, including its extension across the
zeros of theta by the identity theorem, is
\srcid{ThetaUniformisationFrameEqAutomorphyMul.lean}
(\srcid{7e0fe6b56293d92ad64e72bacde8d9aeac55f766}).
The exact box formula is
\srcid{ExistsHolomorphicThetaOfLevelBasis.lean}
(\srcid{fa578ce841cf27bb7ad5bf5c6047b9011a4b49f2}).
Moving zeros is
\srcid{SimpleZerosMoveHolomorphically.lean}
(\srcid{f2458805d37961ce99100633147e62bb53bf388d});
varying contractions is
\srcid{ExistsDifferentiableOnTorsSq.lean}
(\srcid{f93085db0421f99f72c9d2061f17e30b7141acf2}).
Their exact assembly, including the chart-factor cancellation,
is \srcid{KoszulInvariantLocallyHolomorphicPeriodMatrix.lean}
(\srcid{34a2165e6ee4a1b7242efdac1129f8328adcddbe}).
All unqualified filenames after the first one in this paragraph
are in \srcid{ThetaLawProofs/}.

\subsection{The source's weight-zero constancy argument}
\label{sec:ftb-key-weight-zero}

This subsection supplies the weight-zero statement used to normalize
the holomorphic key invariant.  It follows the frozen source's
Fourier estimate and rank-one modular curves.  The sole
one-dimensional modular input is the usual theorem that a
holomorphic modular form of weight zero for
\(\mathrm{SL}_2(\mathbb Z)\), bounded at its cusp, is constant.
It follows, for example, by the removable-singularity theorem and
the maximum principle on the compactified standard fundamental
domain.  No higher-dimensional compactification or boundary
extension theorem is used here.

\begin{srlemma}[Fourier bound on a Siegel cone]
\label{lem:ftb-ph-koecher}
Let \(g\geq2\), \(N\geq1\), and let \(F\) be holomorphic on
\(\mathfrak H_g\).  Suppose
\[
 \begin{aligned}
 F(\tau+NB)&=F(\tau) &&(B=B^t\in M_g(\mathbb Z)),\\
 F(U\tau U^t)&=F(\tau)
 &&(U\in\mathrm{SL}_g(\mathbb Z),\ U\equiv1\bmod N).
 \end{aligned}                                             \tag{W1}\label{eq:ftb-ph-periods}
\]
For every \(c>0\) there is a finite real \(K_c\) such that
\[
       |F(\tau)|\leq K_c\quad
          \text{whenever }\operatorname{Im}\tau\geq c1.
 \tag{W2}\label{eq:ftb-ph-cone-bound}
\]
\end{srlemma}
\begin{proof}
For clarity we give the source's full-matrix Fourier construction;
this avoids off-diagonal half-integer conventions.
Put
\[
 \Phi(z)=N(z+z^t),\qquad
 \mathcal U=\{z\in M_g(\mathbb C):\operatorname{Im}\Phi(z)>0\},
 \qquad f(z)=F(\Phi(z)).
\]
This is an open convex tube.  The function \(f\) has period \(1\)
in each of its \(g^2\) coordinates, by \eqref{eq:ftb-ph-periods}.
For \(m\in M_g(\mathbb Z)\), the integral
\[
 c_m=\int_{[0,1]^{g^2}}
 f(x+z)e^{-2\pi i\langle m,x+z\rangle}\,dx
 \tag{W3}\label{eq:ftb-tube-coeff}
\]
is independent of \(z\in\mathcal U\).
Indeed the integrand is holomorphic and has unit periods.
Its cube integral is holomorphic in \(z\), by differentiation
under an integral over a compact cube, and invariant under every
real translation, by the periodic cube-change formula.
Its real derivatives vanish; complex differentiability therefore
makes all its complex derivatives vanish.  Convexity makes it
constant.  Here \(\langle m,z\rangle=\sum_{ij}m_{ij}z_{ij}\).

Set \(a_m=|c_m|\).
If a real imaginary slice is bounded by \(M\), its Fourier
coefficient estimate gives
\[
          a_m e^{-2\pi\langle m,Y\rangle}\leq M.
 \tag{W4}\label{eq:ftb-tube-coeff-bound}
\]
At a fixed \(z\), any summable nonnegative majorant \(b_m\) of
the left side gives \(|f(z)|\leq\sum_m b_m\).
To justify this last assertion without assuming pointwise Fourier
convergence, regard \(x\mapsto f(x+z)\) as a continuous function
on the unit real torus.  Its Fourier coefficients are the
corresponding expressions in \eqref{eq:ftb-tube-coeff}.
An absolutely summable coefficient sequence has a uniformly
convergent Fourier series.  Fourier uniqueness, or the elementary
approximate-identity argument on that torus, identifies its sum
with the original continuous function.  Taking norms proves the
assertion.

The unit periods make \(F\) bounded on each fixed positive
imaginary slice: reduce the real symmetric entries to a compact
cube.  Fix \(c>0\), and let \(M\geq0\) bound the slice
\((c/2)1\).  Set \(s=c/(4N)>0\).
For every antisymmetric real \(A\), the slice \(Y=s1+A\)
of \(\mathcal U\) has
\(N(Y+Y^t)=(c/2)1\); hence
\[
 a_m e^{-2\pi\langle m,s1+A\rangle}\leq M.
 \tag{W5}\label{eq:ftb-anti-test}
\]
For every \(U\in\mathrm{SL}_g(\mathbb Z)\) congruent to \(1\bmod N\),
invariance and conjugation by the integral inverse of \(U\)
give the same bound on the slice \((c/2)UU^t\).  Thus
\[
 a_m e^{-2\pi s\,\operatorname{tr}(U^tmU)}\leq M.
 \tag{W6}\label{eq:ftb-congruence-test}
\]

If \(a_m>0\), \eqref{eq:ftb-anti-test} forces \(m=m^t\):
otherwise choose \(A=t(E_{ij}-E_{ji})\) and let
\(\langle m,A\rangle\) tend to \(-\infty\).
Next \(m\) is positive semidefinite.
If not, continuity gives a rational vector \(x\) with \(x^tmx<0\);
clear denominators to take \(x\in\mathbb Z^g\).
Choose \(i\) with \(x_i\ne0\), choose \(j\ne i\), and put
\(y=x_je_i-x_ie_j\ne0\).
Then \(y^tx=0\) and
\[
 U_k=1+Nkxy^t\in\mathrm{SL}_g(\mathbb Z),\quad
 U_k\equiv1\bmod N,\quad
 \operatorname{tr}(U_k^tmU_k)
 =\operatorname{tr}m+2Nk\,x^tmy
                 +(Nk)^2(x^tmx)\,y^ty .
\]
The determinant formula for a rank-one perturbation proves the
first assertion.  The last polynomial tends to \(-\infty\).
It contradicts \eqref{eq:ftb-congruence-test}, because \(a_m>0\).
Rational vectors are dense, so this proves real positive
semidefiniteness as well.

For such an \(m\), put \(t=\operatorname{tr}m\geq0\).
The inequalities obtained by evaluating its quadratic form at
\(e_i\pm e_j\) give \(|m_{ij}|\leq t\).  Consequently
\[
       \sum_{ij}|m_{ij}|\leq g^2t,\qquad
       \langle m,\operatorname{Im}\tau\rangle\geq ct
                    \quad(\operatorname{Im}\tau\geq c1).
\]
Use \eqref{eq:ftb-tube-coeff-bound} at \(Y=s1\).
With \(\varepsilon=\pi c/(2Ng^2)>0\) the preceding estimates give
\[
 a_m e^{-(\pi/N)\langle m,\operatorname{Im}\tau\rangle}
       \leq M\exp\!\left(-\varepsilon\sum_{ij}|m_{ij}|\right).
 \tag{W7}\label{eq:ftb-ph-summable}
\]
The same inequality is trivial if \(a_m=0\).
The right side is summable over \(M_g(\mathbb Z)\): it is a
finite product of two-sided geometric series.
Apply the torus Fourier bound at \(z=\tau/(2N)\).
This proves \eqref{eq:ftb-ph-cone-bound}, with exactly the
source's summable majorant.
\end{proof}

\begin{srlemma}[A finite modular family is constant]
\label{lem:ftb-finite-elliptic-family}
Let \(\mathcal F\) be a finite set of holomorphic functions on
the upper half-plane, each bounded as the imaginary part tends
to infinity.  Suppose precomposition by \(S:z\mapsto-1/z\)
and \(T:z\mapsto z+1\) takes \(\mathcal F\) into itself.
Then every function in \(\mathcal F\) is constant.
\end{srlemma}
\begin{proof}
Precomposition is injective, since \(S,T\) are bijections, so it
permutes the finite set.
For each \(w\in\mathbb C\),
\(\prod_{f\in\mathcal F}(w-f(z))\) is therefore a holomorphic
weight-zero modular function for the full modular group, bounded
at infinity and hence at all its equivalent cusps.
The one-dimensional modular input makes it constant.
Fix \(z_0\).  Taking \(w=f(z)\) in that polynomial identity shows
that \(f(z)\) belongs to the finite set
\(\{h(z_0):h\in\mathcal F\}\).
The upper half-plane is connected; continuity makes \(f\)
constant.
\end{proof}

\begin{srproposition}[Holomorphic symplectic invariants have weight-zero constancy]
\label{prop:ftb-ph-holo}
If \(g\geq2\) and \(H:\mathfrak H_g\to\mathbb C\) is holomorphic
and invariant under \(\mathrm{Sp}_{2g}(\mathbb Z)\), then \(H\)
is constant.
\end{srproposition}
\begin{proof}
The symplectic translations and diagonal symplectic matrices give
\eqref{eq:ftb-ph-periods} with \(N=1\).
Lemma~\ref{lem:ftb-ph-koecher} gives a bound on every cone
\(\operatorname{Im}\tau\geq c1\).
We use this bound on the following explicit rank-one curves.

Fix a coordinate \(e\), write \(\epsilon=e_e\), and suppose
\(Z=Z^t\) has zero \(e\)-th column and
\[
       \operatorname{Im}Z+yuu^t>0
                    \quad\text{for all }u_e=1,\ y>0.
\]
For real symmetric \(R\) with \(R_{ee}=0\), and real \(u_e=1\),
put \(C_{R,u}(z)=Z+R+zuu^t\).
For a fixed integer \(N\geq1\), consider all functions
\(z\mapsto H(C_{R,u}(z))\) satisfying
\[
    Nu_i\in\mathbb Z,\quad NR_{ie}\in\mathbb Z,\quad
                             N^2R_{ij}\in\mathbb Z.
 \tag{W8}\label{eq:ftb-curve-denominators}
\]
There are only finitely many such \emph{functions}.
To see this, set \(k_e=0,\ k_i=\lfloor u_i\rfloor\) for \(i\ne e\).
Conjugation by the integral unimodular rank-one shear
\(1-k\epsilon^t\) changes the curve parameters to
\[
 u_1=u-k,\qquad
 R_1=R-k(Re)^t-(Re)k^t .
\]
It fixes \(Z\), since the \(e\)-th row and column of \(Z\) are zero.
Subtract the symmetric integral matrix of floors of \(R_1\);
symplectic translation changes no value of \(H\).
Now \(u_1\) and the fractional entries of \(R_1\) lie in
\([-1,1]\), and their denominators divide \(N^2\).
There are finitely many choices.

Each function is holomorphic and bounded at infinity.
For the bound, its imaginary part for \(y\geq1\) is at least
\(\operatorname{Im}Z+uu^t>0\); choose \(c1\) below this fixed
positive matrix and apply the cone bound.
The finite family is stable under \(T\), which changes the
parameters, up to the integral \(\epsilon\epsilon^t\) translation,
to
\[
                         (R+uu^t-\epsilon\epsilon^t,\ u).
\]
It is also stable under \(S\).  The symplectic partial inversion
in coordinate \(e\) gives the parameters
\[
 u'=\epsilon+Re,\qquad
 R'=R-u(\epsilon+Re)^t-(\epsilon+Re)u^t
                                      +2\epsilon\epsilon^t.
 \tag{W9}\label{eq:ftb-curve-S}
\]
These have \(u'_e=1,\ R'_{ee}=0\), and satisfy
\eqref{eq:ftb-curve-denominators}: multiply each formula by \(N\)
or \(N^2\) and use the displayed integrality conditions.
The partial-inversion identity is checked entry by entry after
using \(Z e=0,\ u_e=1,\ R_{ee}=0\); it asserts that
\(H(C_{R,u}(-1/z))=H(C_{R',u'}(z))\).
Lemma~\ref{lem:ftb-finite-elliptic-family} makes every such
rational-parameter curve constant.

For real \(u_e=1\) and \(v_e=0\), take
\(R=\epsilon v^t+v\epsilon^t\).
Approximate every coordinate of \(u,v\) by
\(\lfloor(n+1)x\rfloor/(n+1)\).
The \(e\)-coordinates remain exactly \(1,0\), respectively, and
the denominator conditions hold with \(N=n+1\).
Passing to the limit by continuity gives constancy on this
real-parameter curve as well.

We now exhibit such a curve through every \(\tau=X+iY\).
Set
\[
 h=Y_{ee},\quad u=Ye/h,\quad
 v=Xe-X_{ee}u,\quad
 R=\epsilon v^t+v\epsilon^t,\quad
 Z=\tau-R-\tau_{ee}uu^t .
\]
Then \(h>0,\ u_e=1,\ v_e=0,\ Ze=0\) and
\(\operatorname{Im}Z=Y-(Ye)(Ye)^t/h\).
For every \(u'_e=1,\ y>0\),
\[
                   Y-(Ye)(Ye)^t/h+yu'(u')^t>0.
\]
Indeed for \(x\ne0\) put
\(x'=x-(x^tYe/h)e\).
Its quadratic value is
\((x')^tYx'+y(x^tu')^2\).
If \(x'\ne0\), the first term is positive.
If \(x'=0\), \(x\) is a nonzero multiple of \(e\) and the second
term is positive.  Thus the curve has the required domain and
passes through \(\tau\) at \(z=\tau_{ee}\).
Its constancy gives
\[
                     dH_\tau((Ye)(Ye)^t)=0.
\]
Conjugate this assertion by integral unimodular matrices.
Their \(e\)-th rows range over the primitive integer vectors,
so it gives
\(dH_\tau((Yv)(Yv)^t)=0\) for every such \(v\).
The vectors \(e_i\) and \(e_i+e_j\), \(i\ne j\), are primitive
rows of unimodular matrices; their outer squares span all complex
symmetric matrices, by subtracting the diagonal squares.
Since \(Y\) is invertible, \(dH_\tau\) vanishes on the whole
symmetric tangent space.
Finally \(\mathfrak H_g\) is convex.  Integrating along its line
segments proves that \(H\) is constant.
\end{proof}

\begin{srcorollary}[Finite symplectic families]
\label{cor:ftb-finite-sp-family}
Let \(g\geq2\).  A finite collection of holomorphic functions on
\(\mathfrak H_g\), permuted by precomposition by every integral
symplectic transformation, consists of constant functions.
\end{srcorollary}
\begin{proof}
Discard duplicate functions.
For each \(w\in\mathbb C\) the polynomial
\(\prod_f(w-f(\tau))\) is holomorphic and symplectic invariant,
because precomposition is an injective map of this finite set
and hence a permutation.
Proposition~\ref{prop:ftb-ph-holo} makes it constant.
Its roots at a fixed period form a finite set containing every
value of every \(f\).  Connectedness and continuity conclude
exactly as in Lemma~\ref{lem:ftb-finite-elliptic-family}.
\end{proof}

\paragraph{Application to the powered key rule.}
In this paragraph assume explicitly \(g\geq2\).
For a field-functorial section-independent rule satisfying
\eqref{eq:ftb-koszul-homog}, the period construction gives one
holomorphic value function \(U_{a,b}\) for each half-characteristic.
The theta transformation formula permutes this finite family.
More precisely, for \(\delta\in\mathrm{Sp}_{2g}(\mathbb Z)\),
write \(j=C\tau+D\).  The theta formula gives a nonzero multiplier
\(r\) with \(r^8=(\det j)^4\); reparametrizing the uniformization
changes the normal-form data to
\[
        \tau'=\delta\tau,\qquad M'=Mj^{-1},\qquad c'=c/r.
\]
It preserves
\[
                         \frac{(\det M')^4}{(c')^8}
                             =\frac{(\det M)^4}{c^8}.
 \tag{W10}\label{eq:ftb-normal-ratio-sp}
\]
The Gaussian exponential in the theta formula is incorporated
in the reparametrized line frame, not discarded.
The homogeneity law and \eqref{eq:ftb-normal-ratio-sp} show that
the normal value functions are permuted, with no scalar factor.
Apply Corollary~\ref{cor:ftb-finite-sp-family}.
For the zero characteristic there is a constant \(a_0\) such that
\[
             \kappa(x)=a_0
                    \left(\frac{(\det M)^4}{c_0^8}\right)^c
 \tag{W11}\label{eq:ftb-key-unscaled-constant}
\]
whenever \(x\) is in zero-characteristic normal form
\((\tau,M,c_0)\).
To pass from the normalized \(M=1,c_0=1\) case to this formula,
replace the theta section by \(c_0^{-1}\vartheta\) and the Lie
basis by \(M^{-1}v\), then apply
\eqref{eq:ftb-koszul-homog}.
The symplectic transformation formula here is the classical
Poisson theta transformation with its eighth-power identity;
its source implementation uses the squared formula and the
identity theorem to choose \(r\), and then changes the actual
uniformization frame.

This proves constancy, not yet the rationality or integral unit
property of \(a_0\), and not yet an exponent-one key rule.
Those are supplied by the subsequent arithmetic normalization
and connected-family descent.

\paragraph{Frozen source concordance.}
All proof paths below have prefix
\srcid{src/Development/AbelianVarietyTheory/Proofs/ThetaLawProofs/}.
The tube Fourier construction and cone estimate are
\srcid{NormLeTsumOfCubeFourierCoeffLe.lean}
(\srcid{958fca26bfcc625698732b1a0dafee0d51889f7b});
its explicit congruence transvection is
\srcid{ExistsNatMulQuadraticLt.lean}
(\srcid{f07ff82f9eaae326e1ce8ca8a834d8f0c40febf0}).
The finite elliptic family, rational rank-one curves, partial
inversion, rational approximation, positive Schur complement,
and spanning of derivatives are
\srcid{FiniteAbsLeOneMulMemRange.lean}
(\srcid{a525213fbc51d8331f5cf36529a36e5f4385887a}).
The final polynomial argument for a finite symplectic family is
\srcid{ProdCompEqOfFiniteFamily.lean}
(\srcid{0d643fd13242090c4d63d7543da06e3fab784482}).
The exact change of theta normal form is
\srcid{ThetaLawSymplecticNormalForm.lean}
(\srcid{6950e84acb9e8fa4fc102ec0cb3baef827e74e78}),
and the rescaling and finite-characteristic-family application
are \srcid{ExistsRescale.lean}
(\srcid{193bbb1a07dfd9582c3d8692d5710c88dabb4500}).
The realization of a testing key datum at each period is
\srcid{ExistsKeyInNormalAt.lean}
(\srcid{189fd3b621d2d5c557c11140849fa5f549a6b81c}),
with its explicit \(\mathrm{TorusAlgebraic}\) hypothesis.


\subsection{Arithmetic normalization and a powered rule at every prime}
\label{sec:ftb-powered-key-all-primes}

The algebraic inputs here are projectivity and finite presentation of
the polarized abelian scheme; faithful-flat descent of its morphisms,
line bundles and identities; the origin-normalized theorem of the cube;
base change of its positive line-power section modules; and finite
etaleness of multiplication by an invertible integer.  These are the
same algebraic inputs as in the integral Koszul construction.  We keep
their applications explicit.

\begin{srlemma}[Countable descent preserving a local principal ideal base]
\label{lem:ftb-key-countable-descent}
Let \(R\) be a local principal ideal domain and \(x/R\) a key datum.
Every countable \(S\subset R\) is contained in a countable local
principal ideal subring \(R_0\subset R\), local and faithfully flat
in \(R\), over which \(x\) descends as a key datum.
Over a characteristic-zero field it descends to a countable subfield.
\end{srlemma}
\begin{proof}
The source constructs subrings as follows.
For countable \(S'\subset R\), put \(K=\operatorname{Frac}R\), let
\(K_0\) be the subfield generated by \(S'\), and set \(R_0=R\cap K_0\).
It is countable.  Divisibility descends: if \(b=ac\) with
\(a,b\in R_0\), then for \(a\ne0\), \(c=b/a\in R\cap K_0\);
the case \(a=0\) is immediate.  Units of \(R\) lying in \(R_0\)
therefore have their inverses in \(R_0\).
The subring is local and its map to \(R\) is local.
For an ideal \(I\subset R_0\), choose \(a\in I\) with \(aR\)
maximal among \(\{bR:b\in I\}\), using Noetherianity of \(R\).
Divisibility is total in \(R\).  For every \(b\in I\), maximality
then gives \(a\mid b\) in \(R\), hence in \(R_0\).
Thus \(I=aR_0\).  Torsion freeness over this principal ideal domain
makes \(R\) flat over \(R_0\); locality makes it faithfully flat.

Descend a projective presentation and its ample bundle using their
finite equations and gluing charts.  Descend the identity, inverse,
multiplication, and theta line bundle using finitely many further
coefficients.  Enlarge to a subring of the preceding kind.
The product used for multiplication is the descended product.
Faithful flatness tests its group identities and descends the
proper smooth polarized scheme, relative dimension, rigidified
symmetric ample line and principal polarization.
Thus the descended object is a polarized abelian scheme.

Over this local principal ideal ring, choose frames of the line at
the origin, the rank-one theta-section module and the rank-\(g\)
Lie module.  Their pullbacks differ from the specified data by an
actual line-bundle identification, a unit scalar and an invertible
matrix.  Enlarge \(S'\) to include these coefficients and the given
\(S\); the subring construction also includes the inverses of their
unit determinants.  Rescaling the frames gives the original datum
under the source's \(\mathrm{KeyPB}\), including its dual-number
Lie maps.  If \(R\) is a field, locality reflects every nonzero
element of \(R_0\) as a unit, so \(R_0\) is a subfield.
\end{proof}

\begin{srlemma}[Unit realization over a local DVR extension]
\label{lem:ftb-key-dvr-realization}
Fix admissible \(n,l,l',\alpha\), with
\(D=\prod_{j=0}^g(l-jn)(l'-jn)\).
For a key datum over a DVR \(R\) in which \(D\) is invertible,
there is an injective local map \(R\to R'\) to a DVR where
the exact \(\operatorname{KaVal}\) formula
\eqref{eq:ftb-kaform} is realized by a unit.
Its value commutes with field base change of that realization.
\end{srlemma}
\begin{proof}
Choose a basis of the level-\(n\) section module.
Over an algebraic closure of a rational-function extension of the
residue field choose \(g\) coefficient rows algebraically independent
over that field.  The preceding generic-section argument makes
their common zeros transverse.
Lift these finitely many coefficients after a local DVR extension
and completion.  Transcendental residue coefficients are obtained
by \(R[T_1,\ldots,T_b]_{(\pi)}\); further dominating DVR extensions
and completion permit the finitely many additional residue
coordinates.  This is the source's
\srcid{exists\_complete\_discreteValuationRing\_with\_residues}.
No quantitative extension degree is required.

Enlarge again to define the residue common zeros and the
\(D\)-torsion.  Multiplication by \(D\) is finite etale, and the
transverse zero scheme is etale at its zeros.  Hensel lifting over
the complete DVR supplies their sections with distinct reductions;
properness and the fixed zero count give all generic zeros.
The actual split-torsion and split-transverse-zero hypotheses
of Proposition~\ref{prop:ftb-integral-koszul-invariant} now hold.
It gives the unit and its termwise base-change identities.
The extension supplies a realization; it does not change the
section-independent field value.
\end{proof}

\begin{srproposition}[Rules away from one integer and field descent]
\label{prop:ftb-key-away-rule}
For each \(g\geq1\) and prime \(p\), there are positive integers
\(c,M\), \(p\nmid M\), and a section-independent rule on
characteristic-zero fields that is functorial for key base change,
satisfies
\[
                 a^{8c}\kappa_{\rm new}
                          =(\det B)^{4c}\kappa_{\rm old},
\]
has holomorphic normalized-period values for every
half-characteristic, and is a unit on each DVR where \(M\) is a unit.
Its unit specializes correctly if the residue field has
characteristic zero.
\end{srproposition}
\begin{proof}
Take \(n=4p,\ l=gn+1,\ l'=gn+1+p,\ M=D\);
Lemma~\ref{lem:ftb-key-numerical-data} gives \(p\nmid M\).
Proposition~\ref{prop:ftb-key-tuple-independent} defines a unique
nonzero value over algebraically closed fields and every generic
coefficient extension, agrees with every admissible tuple, and
gives field functoriality.  Its homogeneity is
\eqref{eq:ftb-koszul-homog}.
The preceding DVR lemma supplies a unit realization after a
local extension.  For a complex datum, descend to a countable
subfield, choose finitely many algebraically independent complex
coefficients, and obtain an actual transverse tuple.
Proposition~\ref{prop:ftb-key-holomorphic-period} supplies and
patches the local holomorphic value functions.

For an arbitrary characteristic-zero field \(K\), pull the datum
to \(\overline K\).  Every \(K\)-automorphism fixes its value by
functoriality and uniqueness.  An algebraic element with this
property lies in \(K\), since its minimal polynomial is separable.
This gives the unique descended nonzero value and its functoriality
under every field map.  To descend the homogeneity identity,
base change the two specified frames simultaneously and use
injectivity and the determinant base-change formula.

Let \(z\in\operatorname{Frac}(R)^\times\) be the descended value,
and \(u\in(R')^\times\) its unit realization.  In a common
algebraically closed fraction field their images agree.
Either \(z\in R\) or \(z^{-1}\in R\).
In the first case its image is \(u\), so the local map \(R\to R'\)
reflects it as a unit.  The second case applies the same argument
to \(z^{-1}\).  This descends the unit without requiring a finite
extension degree.  At a characteristic-zero residue map, pass to
a common algebraic closure of the residue fields of \(R'\) and
\(R\).  The actual determinant, cubical and Koszul formula
specializes there to the residue rule.  Locality makes the
residue maps factor through the original residue field; its
injective embedding in that closure proves the required
specialization equality.
\end{proof}

\begin{srlemma}[The precisely normalized elliptic test]
\label{lem:ftb-key-elliptic-discriminant}
For a smooth Weierstrass curve \(W/R\), the origin line bundle,
its canonical section and rigidification, and the tangent
\([\epsilon:-1:0]\), form a one-dimensional key datum.
For every complex normal form of its base change,
\[
                         \frac{M^4}{c_0^8}=s(\Delta_W).
 \tag{A1}\label{eq:ftb-key-elliptic}
\]
\end{srlemma}
\begin{proof}
The origin parameter on the dual-number tangent
\([\epsilon:-1:0]\) has value \(-\epsilon\).
The rigidified frame has value \(1\) at the origin, so the theta
function vanishes there and its derivative in that tangent
direction is \(-1\).
The three even genus-one half-characteristics give even theta
functions, by the substitution \(h\mapsto-h-2a\) in the series.
Their derivative is zero.  Hence the characteristic is the odd
one \((\tfrac12,\tfrac12)\).

Put \(b=2\pi i/M\).  The actual period lattice on that tangent line
is \(b(\mathbb Z+\tau\mathbb Z)\), and the derivative assertion says
\(\theta'[\tfrac12,\tfrac12](0)=-b/c_0\).
Jacobi's derivative formula in this normalization and the
Weierstrass period formula are
\[
 \theta'[\tfrac12,\tfrac12](0)^8=(2\pi)^8\Delta_{\rm mod}(\tau),
 \qquad
 \Delta\bigl(b(\mathbb Z+\tau\mathbb Z)\bigr)
                       =(2\pi/b)^{12}\Delta_{\rm mod}(\tau).
\]
Substitution gives the latter discriminant as
\((2\pi)^4b^{-4}c_0^{-8}=M^4/c_0^8\), since \(i^4=1\).
The actual Weierstrass uniformization identifies this lattice
discriminant with \(s(\Delta_W)\): its coordinate functions
satisfy the equation of \(W\), with the fixed origin parameter
\(t(0)=0,t'(0)=1\) and its prescribed leading coordinate limits.
These determine the same Weierstrass \(g_2,g_3\), equivalently
the standard periodic-solution uniqueness argument.
Thus both the tangent scale and the discriminant scale are fixed.
The analytic inputs are exactly the Jacobi derivative and
Weierstrass lattice identities displayed here.
\end{proof}

\begin{srproposition}[A normalized powered rule at every prime]
\label{prop:ftb-normalized-powered-key}
For each \(g\geq1\), a positive integer \(N(g)\) and a field-functorial
rule \(\rho_g\) exist with
\[
 \rho_g(x)=
       \left(\frac{(\det M)^4}{c_0^8}\right)^{N(g)}
 \tag{A2}\label{eq:ftb-powered-normal-value}
\]
in every complex half-characteristic normal form.
It obeys
\(a^{8N(g)}\rho_g(x_{\rm new})
=(\det B)^{4N(g)}\rho_g(x_{\rm old})\).
For every characteristic-zero DVR key datum its generic value
is a unit, with the correct specialization whenever the residue
field has characteristic zero.
\end{srproposition}
\begin{proof}
Choose an away-from-\(M\) rule in dimension \(2g\), exponent \(c>0\).
The finite-family argument gives the zero-characteristic constant
\(a_0\) of \eqref{eq:ftb-key-unscaled-constant}.
Evaluate this rule on \(x^2\).
Its doubled half-characteristic is even, and an integral symplectic
transformation takes it to zero.  The transformation preserves
\((\det M)^4/c_0^8\) by
\eqref{eq:ftb-normal-ratio-sp}; the square datum squares that ratio.
The resulting rule in dimension \(g\), exponent \(n=2c\), therefore
has normal value \(a_0((\det M)^4/c_0^8)^n\) for every
half-characteristic.  The field and DVR laws pass to the square.

For any integral Weierstrass curve \(W\), take its key datum over
\(\mathbb Z[\Delta_W^{-1}]\) and then its \(g\)-th power.
Lemma~\ref{lem:ftb-key-elliptic-discriminant} makes the value
\(a_0\Delta_W^{gn}\).  The field rule is defined over \(\mathbb Q\);
dividing its rational value by \(\Delta_W^{gn}\) proves
\(a_0=r\in\mathbb Q^\times\).
Whenever \(M\) and \(\Delta_W\) are DVR units, both the rule and
this discriminant power are units, so \(r\) is a unit there.
Use the two integral curves \(W_{j=0}\), with discriminant \(-27\),
and \(W_{j=1728}\), with discriminant \(-64\).
Both give the same \(r\), because they give the same complex \(a_0\).
In a local ring either \(3\) or \(1-3=-2\) is a unit.
One of these two discriminants is therefore a unit.
Thus \(r\) is a unit at every DVR where \(M\) is a unit.
Dividing by \(r\) gives exact normal value and preserves unit
and characteristic-zero specialization laws.

Do this separately for each prime \(p\), using \(M_p\) prime to
\(p\), and denote the normalized rule and exponent by
\((\rho_p,n_p)\).
For any two,
\[
                             \rho_p^{n_q}=\rho_q^{n_p}
 \tag{A3}\label{eq:ftb-powered-cross}
\]
on every characteristic-zero field.  Descend a datum to a
countable subfield and embed it in \(\mathbb C\); the equality
there is immediate from the exact normal values.
Functoriality and injectivity descend it.

Fix the rule at \(p=2\), setting \(N(g)=n_2\).
At a DVR where \(M_2\) is invertible it is already a unit.
Otherwise the residue field has a prime characteristic \(p\)
dividing \(M_2\).  The rule \(\rho_p\) is a unit there since
\(p\nmid M_p\); \eqref{eq:ftb-powered-cross} makes
\(\rho_2^{n_p}\) a unit.  Its valuation is
\(n_pv(\rho_2)\); positivity of \(n_p\) forces
\(v(\rho_2)=0\).
A characteristic-zero residue field makes every positive integer,
including \(M_2\), a unit; its specialization is therefore always
covered by the first case.  No excluded residue prime survives.
\end{proof}

\paragraph{Frozen source concordance and remaining producer.}
All paths have prefix
\srcid{src/Development/AbelianVarietyTheory/Proofs/ThetaLawProofs/}.
The local principal ideal subring is
\srcid{ExistsCountableLocalPidSubring.lean}
(\srcid{80f4e931782aab0999790894204dc2b851005194});
the actual scheme/group/line/frame descent is
\srcid{KeyDataDescendCountableSubrings.lean}
(\srcid{d06ec6cf161d65aff27cc78f9f3693252a1b58d2}).
DVR lifting and the away rule are
\srcid{ExistsKeyInvariantAfterDvrExtension.lean}
(\srcid{e9d5f9addb9796136f989270119a664641668639})
and \srcid{ExistsIntegralKeyInvariantRuleCharZero.lean}
(\srcid{86b60fd21a1e75dd48b4971ad05dc6e34e3f4228}).
Field descent and reflection of units are
\srcid{ExistsKeyPBPair.lean}
(\srcid{db42d4b4b8b9d8f9cd9a17237b338a72ed354076}).
The elliptic test is
\srcid{ExistsEllipticKeyDataValueEqDiscriminant.lean}
(\srcid{0e56730263497260189223973dbdaec4d38d1b09}).
Squaring, the rational tests and division by their common unit are
\srcid{ExistsNormalisedKeyInvariantRuleCharZero.lean}
(\srcid{e9c21ea730a06bb841cbfc137a7964f8f1c461c2}).
Cross powers and removal of all bad primes are
\srcid{ThetaKeyRuleExistsOfAlgebraicInvariants.lean}
(\srcid{24e1e9047fad6a46b1a8b2f4ecb81f4bbd835610}).

This is a powered rule with \(N(g)>0\), not yet the exponent-one
rule required by the key section.  Taking an arbitrary root would
lose field functoriality and specialization.  The source next fixes
that root by connected algebraic families with a known arithmetic
point.  The complex uniformization and polarized-torus
algebraization interfaces retained in the preceding subsection
also remain in force.


\subsection{The chosen root in a connected family}
\label{sec:ftb-key-connected-root}

For a complex key datum in a normal form \((\tau,M,a,b,c)\), put
\[
                            t(x)=\frac{(\det M)^4}{c^8}.
 \tag{R1}\label{eq:ftb-key-root-value}
\]
We first justify that this is a value of the datum, not of the
chosen normal form.

\begin{srlemma}[Independence of the normal form]
\label{lem:ftb-key-normal-value-unique}
With the actual complex theta uniformization and the Poisson theta
transformation law, \eqref{eq:ftb-key-root-value} is independent of
the half-characteristic normal form.
\end{srlemma}
\begin{proof}
Two exponential uniformizations of the same complex abelian variety
differ by a complex linear map \(\phi\); their specified tangent
isomorphisms determine it.  It maps their period lattices to each
other.  The quotient of their local line frames glues, by the
identical transition functions, to a nowhere-zero entire \(e(z)\).
Both use the same origin rigidification, so \(e(0)=1\).
The theta identity has the form
\[
        c'\theta[a',b'](z,\tau')
                 =e(z)c\theta[a,b](\phi z,\tau).
 \tag{R2}\label{eq:ftb-normal-frame-compare}
\]
The integer lattice map is symplectic.  Here is the actual
automorphy check, which avoids assuming that a group isomorphism
preserves an unspecified polarization.  Write the images of its
\(2g\) period generators as \(m_s+\tau n_s\).
The two automorphy laws and \eqref{eq:ftb-normal-frame-compare}
give the difference equation for the nowhere-zero entire frame
quotient.  Taking a global holomorphic logarithm on \(\mathbb C^g\)
gives
\(q(z+p_s)-q(z)=\ell_s(z)+d_s\), with linear \(\ell_s\);
the possible \(2\pi i\) multiples are absorbed in \(d_s\).
Commuting two differences yields
\(\ell_s(p_t)=\ell_t(p_s)\).
Substitute the linear terms \(-2\pi i n_s\cdot\phi z\)
and their primed counterparts from \eqref{eq:ftb-level-aut}.
The quadratic \(\tau\) terms cancel by symmetry, giving equality
of the two integral alternating forms
\(n_s\cdot m_t-n_t\cdot m_s\).
Thus the integer lattice matrix is symplectic.

Consequently \(\tau'=\delta\tau\) and
\(\phi z=j^tz\), \(j=C\tau+D\), for that symplectic matrix.
The theta transformation gives the multiplier \(r\) with
\(r^8=(\det j)^4\).
After its Gaussian exponential is absorbed, compare the two
theta functions at the same period.  The entire quotient \(\varrho\) between these same-period
functions has sign multipliers on each unit and \(\tau\)-period,
because both characteristics are half-integral.
Thus \(\varrho^2\) has every period of the full lattice.
It is bounded on a compact real period parallelepiped, hence
bounded on \(\mathbb C^g\), and is constant by Liouville.
Since \(\varrho\) is nowhere zero, differentiating its constant
square shows that \(\varrho\) itself is constant.
All its sign multipliers are therefore \(1\).
The two characteristics agree modulo integers; our representatives
\(0,\tfrac12\) make them equal.  Its value at zero then gives
\(c^8=(c')^8(\det j)^4\).
The tangent comparison gives \(M=M'j\).
Taking determinants proves
\((\det M)^4/c^8=(\det M')^4/(c')^8\).
\end{proof}

We will use the following analytic family input, with its exact
scope.  On a small ball of parameters for a smooth proper polarized
abelian scheme, its relative exponential and the chosen
dual-number tangent frame admit jointly holomorphic local
coordinates; a fixed integral symplectic basis of the period
lattice continues holomorphically; its torsion sections and Weil
pairing continue with the same labels.  The resulting normalized
period has positive imaginary part, and the theta section has,
on each line chart, its actual local nowhere-zero quotient by
one of the finitely many half-characteristic theta functions.
These are respectively the source's relative exponential,
period-continuation, torsion/Weil, and normal-fibre statements.
We prove the finite-choice and invariant-continuity step below.
Construction of this family input is distinct from the fixed
period holomorphy in
Proposition~\ref{prop:ftb-key-holomorphic-period}.

\begin{srlemma}[The source's circle quotient argument]
\label{lem:ftb-circle-quotient}
Let \(f(b,z),\theta(b,z)\) be holomorphic on an open subset of
finite-dimensional complex parameter and fibre spaces.
Suppose \(\theta(b_0,-)\) is not the zero germ at \(z_0\).
If, for parameters in a set \(S\), the fibre function \(f/\theta\)
extends holomorphically, then near \((b_0,z_0)\) there is a
holomorphic \(G(b,z)\) satisfying \(G\theta=f\) on parameters
in \(\overline S\).
If a quotient \(F\theta=f\) is defined on all fibres and is
holomorphic in each fibre, then it is jointly holomorphic.
\end{srlemma}
\begin{proof}
Choose a complex direction \(v\) on which the germ of
\(\theta(b_0,z_0+tv)\) is not identically zero.
Such a direction can also be obtained by joining \(z_0\) to a
nearby point where \(\theta\ne0\).
Its zeros on that line are isolated, so choose a sufficiently
small \(r>0\) for which theta is nonzero on \(|t|=r\).
Compactness of the circle and continuity keep it nonzero on
that same circle for all \((b,z)\) near \((b_0,z_0)\), and
keep the whole closed disk inside the chart.
Set
\[
 G(b,z)=\frac1{2\pi}\int_0^{2\pi}
       \frac{f(b,z+re^{iu}v)}{\theta(b,z+re^{iu}v)}\,du .
 \tag{R3}\label{eq:ftb-circle-average}
\]
Its denominator is bounded away from zero locally on the
compact circle.  Differentiating under that compact integral
makes \(G\) jointly holomorphic.
For \(b\in S\), the holomorphic fibre quotient has its value at
\(z\) equal to this circle average, by the one-variable
holomorphic mean-value formula.  Thus \(G\theta=f\) there.
Continuity, applied at each fixed \(z\), extends this equality
to \(b\in\overline S\).
In the second assertion the same circle formula equals \(F\)
on every nearby fibre, proving joint holomorphy.
No extension by a generic-fibre value is used.
\end{proof}

\begin{srlemma}[Finite theta choices and the rigidified scalar]
\label{lem:ftb-finite-frame-continuity}
Assume the actual local family input just stated.
The half-characteristic occurring in its normalized theta frames
is locally constant, and the scalar \(c(b)\) fixed by the
origin rigidification is continuous.
Consequently \(t(x_b)\) is continuous.
\end{srlemma}
\begin{proof}
In each line chart let \(f_j(b,z)\) be the actual algebraic
theta section read on the relative exponential.
It is holomorphic in both variables.
For each fibre, its quotient by the selected theta function
is locally a holomorphic unit.  Nonzero theta is nonzero on
a dense subset of each fibre by the identity theorem, so
these local quotients glue uniquely.

Suppose a different member \(k\) of the finite
half-characteristic set occurs at parameters approaching
\(b_0\).  Apply Lemma~\ref{lem:ftb-circle-quotient} to
\(f_j/\theta_k\), and also to its inverse
\(\theta_k/f_j\), along these parameters.
At the limiting fibre the two holomorphic quotients multiply
to \(1\): this equality holds on the dense set where the
functions are nonzero and then everywhere by continuity.
Compare with the unit quotient for the characteristic selected
at \(b_0\).  On overlapping line charts these quotients glue,
since their transition functions agree; dense nonvanishing
makes the gluing unique even at theta zeros.
They give a nowhere-zero entire ratio between theta functions
of characteristics \(k\) and \(k_0\).
The characteristic uniqueness in the preceding normal-form
argument forces \(k=k_0\), a contradiction.
The fibres of this finite characteristic choice are therefore
closed locally, hence open locally, proving local constancy.

With that characteristic fixed,
Lemma~\ref{lem:ftb-circle-quotient} makes the chart quotient
jointly holomorphic.
At \(z=0\) its value is \(c(b)\) times the actual line
rigidification value \(r_j(b)\).
That value is holomorphic and nonzero.
Dividing gives continuity of \(c(b)\), including when the
theta function itself vanishes at the origin.

For completeness, the tangent matrix is continuous as well.
Let \(P(b)\) be the continued first half of the period basis;
its inverse \(h(b)\) is holomorphic by the adjugate formula.
The vector \(W_i(b)\) in exponential coordinates representing
the \(i\)-th dual-number tangent is holomorphic: choose
\(g\) local functions whose derivatives give an invertible
Jacobian at the identity, and solve that finite system.
Both its coefficients and its tangent evaluations are
holomorphic in \(b\).  The normal-form matrix has rows
\(2\pi i\,h(b)W_i(b)\); thus its determinant is continuous.
It is nonzero, and \(c(b)\ne0\).
Formula \eqref{eq:ftb-key-root-value} proves continuity of \(t\).
\end{proof}

\begin{srlemma}[A powered rule on a smooth base is an actual unit]
\label{lem:ftb-powered-rule-smooth-unit}
Let \(k\) have characteristic zero, let \(S\) be a smooth integral
finitely generated \(k\)-algebra, and let \(y/S\) be a key datum.
The powered rule of
Proposition~\ref{prop:ftb-normalized-powered-key} is represented
by a unit \(u\in S^\times\), at every field-valued point of \(S\).
In particular
\[
                         t(y_s)^{N(g)}=s(u)
 \tag{R4}\label{eq:ftb-family-powered-unit}
\]
at all complex points.
\end{srlemma}
\begin{proof}
At every height-one prime the normal ring \(S\) localizes to a DVR.
The generic rule value and its inverse belong to that DVR
by its unit law.  The height-one intersection theorem for a
Noetherian normal domain puts both in \(S\), so they are
mutual inverses there.  This gives \(u\).

We must still identify its specialization in higher codimension.
The source proves this over a regular local ring \(A\) by
induction on its dimension.
Dimension zero is a field.  Otherwise choose
\(\pi\in\mathfrak m_A\setminus\mathfrak m_A^2\).
The quotient \(A/(\pi)\) is a regular local domain of dimension
one less, and \(A_{(\pi)}\) is a DVR.
The rule's characteristic-zero DVR specialization identifies
the generic value on \(A/(\pi)\) with \(u\bmod\pi\).
Apply the induction hypothesis to that quotient and its closed
point.  Every field-valued point with kernel \(\mathfrak m_A\)
factors through its residue field, so field functoriality gives
the assertion at that point.
Localizing \(S\) at any prime reduces its point to this case.
Thus the unit represents the rule at all field-valued points,
not only generically.
The exact complex normal value gives \eqref{eq:ftb-family-powered-unit}.
\end{proof}

\begin{srlemma}[A root cover with two lifted points]
\label{lem:ftb-key-connected-root-cover}
Let \(K\) be algebraically closed of characteristic zero,
\(R/K\) a smooth integral finitely generated algebra, and
\(u\in R^\times\).  Let \(n\geq1\), \(p_0,p:R\to K\), and
\(a_0^n=p_0(u)\).
There is a smooth integral \(K\)-algebra \(B\), a map
\(\phi:R\to B\), a unit \(r\in B^\times\) with
\(r^n=\phi(u)\), and lifts \(q_0,q:B\to K\) of \(p_0,p\),
such that \(q_0(r)=a_0\).
\end{srlemma}
\begin{proof}
Use induction on the prime factors of \(n\).
For a prime \(l\), if \(u=\rho^l\) in \(R\), take \(B=R\)
and replace \(\rho\) by
\((a_0/p_0(\rho))\rho\); its scalar is an \(l\)-th root of
unity and fixes its \(l\)-th power.
Otherwise \(X^l-u\) is irreducible over \(\operatorname{Frac}R\).
Indeed that field contains all \(l\)-th roots of unity; the
conjugates of a root form an orbit of size \(1\) or \(l\).
In the first case the root lies in the fraction field, is
integral over \(R\), and hence lies in \(R\), contrary to the
case assumption.
Take \(B=R[X]/(X^l-u)\).  It is a domain.
Its root is a unit, so its derivative \(lX^{l-1}\) is a unit;
the algebra is finite etale over \(R\), hence smooth over \(K\).
Algebraic closedness gives a lift of \(p\), and the specified
root \(a_0\) gives the desired lift of \(p_0\).
Apply this step and then the inductive root construction for
the remaining prime factors.  The specified root at the first
point is retained at each step.  This is the source's connected
root-cover construction; an arbitrary disconnected root torsor
would not give the required integral base.
\end{proof}

\begin{srlemma}[The root is fixed by one point]
\label{lem:ftb-root-fixed-by-point}
Let \(B\) be a smooth integral finitely generated complex algebra,
\(y/B\) a key datum, \(r\in B^\times\), and \(n\geq1\).
Suppose the actual normal value satisfies
\(t(y_q)^n=q(r)^n\) at every complex point \(q\).
If \(t(y_{q_0})=q_0(r)\) at one point, then
\(t(y_q)=q(r)\) at every point.
\end{srlemma}
\begin{proof}
The complex points of an integral finite-type complex variety
are connected.  One standard proof, used in the source, takes
a Noether-normalization principal open with primitive-element
unramified description, proves that this open is connected
using irreducibility and analytic continuation off its
discriminant, and uses density of that principal open in
complex points.
Smooth algebraic charts are complex analytic charts.
Lemma~\ref{lem:ftb-finite-frame-continuity}, applied in these
charts, makes \(q\mapsto t(y_q)\) continuous.
Evaluation \(q\mapsto q(r)\) is continuous and nonzero.
Their ratio is therefore a continuous map to the finite set
of \(n\)-th roots of unity.  It is constant on the connected
space, and its value at \(q_0\) is \(1\).
\end{proof}

\paragraph{Exact arithmetic-family interface.}
We call a datum \(x/\mathbb C\) \emph{connected to a known point}
if there are a number field \(k\), a smooth geometrically
integral finitely generated \(k\)-algebra \(S\), a key datum
\(y/S\), maps \(s:S\to\mathbb C\) and \(s_0:S\to k\) with
\(s_0\) a \(k\)-map, and \(t_0\in k\), such that
\(x=s^*y\), and for every embedding \(e:k\to\mathbb C\),
\[
                         t((e s_0)^*y)=e(t_0).
 \tag{R5}\label{eq:ftb-connfam}
\]
All identifications here are key pullbacks: line bundle,
rigidification, theta section and ordered Lie frame are included.
Geometric integrality means \(F\otimes_k S\) is a domain for
\emph{every} field extension \(F/k\).
This is the literal source \(\mathrm{ConnFam}\), not just
path connectedness of an analytic moduli space.

\begin{srproposition}[Known families give equivariance under complex field maps]
\label{prop:ftb-key-complex-equivariance}
Assume the local analytic family input above, and that every
complex key datum is connected to a known point in the exact
sense \eqref{eq:ftb-connfam}.
Then for every field map \(f:\mathbb C\to\mathbb C\) and every
key pullback \(x'=f^*x\),
\[
                              t(x')=f(t(x)).
 \tag{R6}\label{eq:ftb-key-complex-funct}
\]
\end{srproposition}
\begin{proof}
For the family in \eqref{eq:ftb-connfam} put
\(e=s|_k\), and base change to
\(R=\mathbb C\otimes_{k,e}S\).
It is smooth integral; \(s\) and \(es_0\) give two complex
points \(p,p_0\).
Lemma~\ref{lem:ftb-powered-rule-smooth-unit} gives a unit
\(u\in S^\times\) with \(t^{N(g)}\) equal to its evaluations.
Apply Lemma~\ref{lem:ftb-key-connected-root-cover} to its
image in \(R\), choosing the root \(e(t_0)\) at \(p_0\).
It gives \(B,r,q,q_0\).
The root-fixed-by-one-point lemma makes \(t(x)=q(r)\).

Now repeat the entire base change by the possibly nonsurjective
embedding \(f:\mathbb C\to\mathbb C\).
The algebra \(\mathbb C\otimes_{\mathbb C,f}B\) is smooth and
integral: a finite-type integral algebra over an algebraically
closed field is geometrically integral.
The points \(fq,fq_0\) lift to this base change, and its root
is the image of \(r\).
At the lifted \(q_0\), \eqref{eq:ftb-connfam} gives the chosen
root \(f(e(t_0))\).
The same root-fixed-by-one-point lemma makes its value at the
lifted \(q\) equal to \(f(q(r))=f(t(x))\).
This is the normal value of \(x'\) by
Lemma~\ref{lem:ftb-key-normal-value-unique}.
No continuity of the field map \(f\) has been assumed or used.
\end{proof}

\paragraph{Frozen source concordance and scope.}
The base prefix is
\srcid{src/Development/AbelianVarietyTheory/Proofs/}.
Normal-value independence is
\srcid{ThetaLawProofs/ThetaFrameValueIndependentNormalForm.lean}
(\srcid{f823bc2749dfa176812b47f1be3902b284c5eada}).
The circle quotient and closure argument are
\srcid{ModuliProofs/DifferentiableOnQuotientOfDifferentiableOnSlice.lean}
(\srcid{efca854eca061af3cefd2f6ce3c9b9f13e6cc61c})
and \srcid{ModuliProofs/EqOfFiniteOfClosedFibres.lean}
(\srcid{7f7959eed9e909e48e267a075fa077ab74d989e9}).
The finite-character and origin-scalar argument is
\srcid{ThetaLawProofs/ExistsCoordinatesOfRelations.lean}
(\srcid{6da8aa085884d27e2caacc57952b26d8dd59704a}).
Its actual family assembly is
\srcid{ThetaLawProofs/ContinuousOnThetaFrameValueHolomorphicFamily.lean}
(\srcid{eae569828b1ee42499288b46d7fcf4dd39dc4062}):
the relative exponential, continued periods, and their
torsion/Weil compatibility remain the explicit family input
stated above.  These are produced in
\srcid{ModuliProofs/ExistsRelativeExpOfParts.lean},
\srcid{LatContMapContinuous.lean},
\srcid{TorsionEqAndWeilIsOfIsHolTo.lean}, and
\srcid{ThetaLawProofs/ExistsNormalOfWeil.lean}.
The smooth-base unit proof is
\srcid{ThetaLawProofs/ExistsUnitThetaKeyRuleSmoothBase.lean}
(\srcid{36413650a428826ff71462b79aded750faaa2d18}).
Connected root constancy is
\srcid{ThetaLawProofs/ThetaFrameValueEqUnitSmoothConnectedBase.lean}
(\srcid{aee5baa0b37f382a5ca6aeb4caca40e115b342aa}).
The integral root covers, geometric base change and equivariance
are \srcid{ThetaLawProofs/IsDomainTensorOfIsAlgClosed.lean}
(\srcid{df53eb7c2f461008fca381d1052984d14cc517d0}).
The arithmetic \(\mathrm{ConnFam}\) interface is expanded here
but its construction is a further producer; it is not inferred
from an arbitrary complex path.


\subsection{The arithmetic family containing the prescribed key datum}
\label{subsec:ftb-arithmetic-key-family}

A key datum \(x=(A,L,\mathrm{rig},\theta,v)\) consists of an abelian
scheme of positive dimension \(g\), a rigidified symmetric principal
line bundle, a fibrewise nonzero section of that line bundle, and
an ordered Lie basis.  A key pullback includes all these data.
Its complex normal value is \(t(x)=\det(M)^4/c^8\), in the
normalization of Lemma~\ref{lem:ftb-key-normal-value-unique}.
Over a number field \(k\), a key datum is \emph{known} if one
\(t_0\in k\) gives its normal value under every embedding of \(k\)
in \(\mathbb C\), and for every corresponding key pullback.
A polarized variety is \emph{potentially known} if, after one
finite extension, every key datum on every further pullback of
that polarized variety is known.
Here the polarization includes the actual line bundle.

We separate two precise moduli inputs from the algebraic family
construction.  The first is a fine principal Siegel family \(F/K\)
at an odd level \(k\ge3\), with \(K\) a number field and
\(\zeta\in K\) a primitive \(k\)-th root.
It is smooth of relative dimension \(g(g+1)/2\), quasiprojective,
and has a universal principal polarization biextension and full
level.  It represents principal Siegel data over reduced bases
with dense characteristic-zero generic fibres, and separates
maps from local affine schemes giving the same Siegel datum.
The second input says that, for
\(\sigma\zeta=\exp(-2\pi i/k)\), any two complex points with
standard period-level data lie on the same connected component
of \(F(\mathbb C)\).
Their actual lower producers are recorded below; neither input
is replaced by an unnamed moduli or connectedness theorem.

\begin{srlemma}[A finite \'{e}tale torsion frame cover]
\label{lem:ftb-torsion-frame-cover}
For a Noetherian ring \(R\), an abelian scheme \(A/R\) of relative
dimension \(g\), and \(n\in R^\times\), there is a finite
\'{e}tale faithfully flat \(R\)-algebra \(R_1\) with
\[
             A[n]_{R_1}\simeq(\mathbb Z/n\mathbb Z)^{2g}_{R_1}.
\]
In particular the labelled sections list the full torsion,
without repetitions, after every algebraically closed field
pullback.
\end{srlemma}
\begin{proof}
The torsion scheme is finite \'{e}tale, and the geometric
\(d\)-torsion has \(d^{2g}\) elements for every \(d\mid n\).
These counts show that each geometric \(n\)-torsion group is
free of rank \(2g\).  Indeed, after choosing \(s<2g\) independent
generators, let \(m\) be the exponent of the quotient by their
span, and write \(n=mq\).  Multiplication by \(m\) has an image
of order \(q^{2g}\).  That image is generated by the images of
the chosen \(s\) generators: \(mx\) belongs to their span by
definition of \(m\), and their independence makes its
coordinates divisible by \(m\).
Thus \(q^{2g}\le q^s\), so \(q=1\).
Lift an element of order \(n\) in the quotient to obtain the
next independent generator.  Induction gives a frame.

In the product of \(2g\) copies of \(A[n]\), remove every
nonzero-coefficient linear-relation locus.  Each such locus
is open and closed, since the diagonal of a finite \'{e}tale
scheme is open and closed.  The remaining subscheme \(W\) is
finite \'{e}tale and is surjective by the geometric argument.
It is affine; let \(R_1\) be its coordinate ring.
Finite flatness and surjectivity give faithful flatness.
The tautological frame is an isomorphism on every geometric
fibre, hence an isomorphism of finite \'{e}tale group schemes.
\end{proof}

\begin{srlemma}[The actual symmetric principal square root]
\label{lem:ftb-symmetric-principal-root-cover}
Let \(R\) be a Noetherian \(\mathbb Q\)-algebra and \(P\) a
principal polarization biextension on an abelian scheme \(A/R\).
Put \(D=\Delta^*P\).
After a finite \'{e}tale faithfully flat change of base there
is a rigidified symmetric relatively ample line \(L\) and an
actual line-bundle isomorphism
\[
                             \lambda_L\simeq P.
 \tag{F1}\label{eq:ftb-family-actual-square-root}
\]
The full \(2\)-torsion is constant on this cover.
\end{srlemma}
\begin{proof}
Apply Lemma~\ref{lem:ftb-torsion-frame-cover} at \(n=2\), and
write \(C=(\mathbb Z/2\mathbb Z)^{2g}\) for the labelled
kernel, with basis \(h_i\).
Biextension additivity and symmetry give, on every test scheme,
\[
 D(x+h)=D(x)+D(h),\quad 2D(h)=0,\quad D(-x)=D(x)
                                                  \quad(2h=0).
 \tag{F2}\label{eq:ftb-diagonal-two-torsion-pic}
\]
Here \(D(x)=x^*[D]\) is a Picard class.
For example, if \(B(x,y)=(x,y)^*[P]\), then \(D(x)=B(x,x)\);
the cross terms are \(2B(x,h)=B(x,2h)=0\).

Choose generators \(\psi_i,\tau_i,j_0\) of, respectively,
\(D-t_{h_i}^*D+D(h_i)\), \(-2D(h_i)\), and \(D-[-1]^*D\).
Their existence is exactly \eqref{eq:ftb-diagonal-two-torsion-pic}.
For \(D^{\otimes2}\), set \(v_i=\psi_i^2\tau_i\) and \(j=j_0^2\).
Then
\begin{align}
 v_i\,t_{h_i}^*v_i&=a_i\in R^\times,\\
 v_i\,t_{h_i}^*v_j&=v_j\,t_{h_j}^*v_i,\\
 v_i\,t_{h_i}^*j&=j\,[-1]^*v_i.                  \tag{F3}
 \label{eq:ftb-diagonal-linearization-relations}
\end{align}
Here is the source's cancellation, rather than an assumption
that arbitrary theta-group lifts commute.
A quotient with zero formal Picard expression is a global unit
on \(A\), hence a base unit by \(\Gamma(A,\mathcal O_A)=R\).
Translations and inversion fix that quotient.
In particular \(\psi_i t_{h_i}^*\psi_i\tau_i\) is fixed by
\(t_{h_j}\), and
\[
 \frac{\psi_i t_{h_i}^*\psi_j}{\psi_j t_{h_j}^*\psi_i}
\]
is fixed by \(t_{h_i}\).
Multiply those two invariance identities and cancel to obtain
the second equality in \eqref{eq:ftb-diagonal-linearization-relations}.
For the third, use the invariance under inversion of
\(\psi_i t_{h_i}^*j_0/(j_0[-1]^*\psi_i)\), and the invariance
under \(t_{h_i}\) of \(j_0[-1]^*j_0\).
Their quotient gives the squared identity with its \(\tau_i\)
factors.  The first equality has zero formal Picard expression.

Adjoin \(b_i^2=a_i\) successively.
Each polynomial \(X^2-a_i\) is monic, its root is a unit, and
its derivative \(2b_i\) is a unit.
These extensions are finite free, \'{e}tale and faithfully flat.
The \(w_i=v_i/b_i\) are commuting involutions.
Their ordered products give generators \(u_c\), for \(c\in C\),
with the actual cocycle identity
\[
                           u_{c+c'}=u_c\,t_c^*u_{c'}.
 \tag{F4}\label{eq:ftb-family-linearization-cocycle}
\]
This follows by putting commuting involutions in basis order.
Their compatibility with \(j\) also persists.

The map \([2]:A\to A\) is finite \'{e}tale and surjective.
Two points with equal doubles differ \'{e}tale locally by a
unique labelled element of its kernel.
Thus \eqref{eq:ftb-family-linearization-cocycle} is a descent
cocycle on the actual fibre product.
The source constructs invariant local generators on each finite
orbit, shrinks the target until they are generators, and
descends their unit ratios.
They glue a line \(M\) with
\[
                                [2]^*M\simeq D^{\otimes2}.
 \tag{F5}\label{eq:ftb-family-double-pullback}
\]
The equality with \(j\) makes the isomorphism with \([-1]^*M\)
invariant under the kernel.  Its ratios descend by faithful
flatness, so \(M\) is symmetric.

To check its polarization, let \(N=\lambda_M-P\) on \(A\times A\).
It is trivial on both axes; the cube law makes it biadditive.
Evaluate \eqref{eq:ftb-family-double-pullback} on test points
to obtain \(N(2x,2y)=0\), and hence \(N(4x,y)=0\).
One must not divide this equality in the Picard group.
Instead use the closed triviality locus for this line family
in the first factor.  It contains the image of \([4]\).
A closed immersion through which a faithfully flat surjection
factors is an isomorphism.
Consequently the line is locally trivial in that factor.
Its trivialization on the other axis and the equality of
global functions with those on the base give a global
trivialization.  This proves \(\lambda_M\simeq P\) as a line
isomorphism.

Replace \(M\) by \(M-e^*M\).
It is rigidified and symmetric, with the same polarization.
The cube law, now using that actual polarization identity,
identifies its square with \(D\).
It is relatively ample because \(D\) is; it is principal
because its polarization is \(P\).
All torsion labels survive these changes of base.
\end{proof}

\begin{srlemma}[Keeping both points on one integral component]
\label{lem:ftb-two-point-etale-integral-component}
Let \(K,K_1\) be number fields with \(K_1\hookrightarrow\mathbb C\).
Let \(R\) be a smooth finitely generated \(K\)-algebra and
\(R\to R_1\) finite \'{e}tale faithfully flat.
Suppose \(p_1:R_1\to\mathbb C\) and \(q:R\to K_1\) are
\(K\)-maps, and \(p=p_1|_R\).
Assume \(p\) and \(q\) agree on all idempotents after extending
scalars to \(\mathbb C\).
There exist a number field \(k\supset K_1\), a smooth
geometrically integral finitely generated \(k\)-algebra \(S\),
a map \(R_1\to S\), a complex point \(s\), and a \(k\)-point
\(s_0\), with
\[
              s|_{R_1}=p_1,\qquad
              s_0|_R=(K_1\hookrightarrow k)q .
 \tag{F6}\label{eq:ftb-two-point-component-maps}
\]
\end{srlemma}
\begin{proof}
Inside \(\mathbb C\), let \(\overline K\) be an algebraic closure
of \(K\).  The ring \(\overline K\otimes_KR_1\) is regular.
Its finitely many irreducible components are disjoint, and
therefore integral open and closed components.
Choose the one containing \(p_1\), with idempotent \(e\).
Its finitely many coefficients descend to a finite extension
\(k_0/K\).
The corresponding localization is geometrically integral:
its extension to \(\overline K\) is a domain, and the
algebraically closed field test gives geometric integrality.

The image of this component in
\(\operatorname{Spec}(\overline K\otimes_KR)\) is open and
closed, by flatness with finite presentation and by finiteness.
It contains \(p\).  Agreement on idempotents puts \(q\) in
that image, so \(q\) has a lift on this same component.
The finitely many coordinates of the lift are algebraic over
\(K_1\), since \(R_1/R\) is finite.
Adjoin them and \(k_0\) to \(K_1\), obtaining \(k\).
The base change of the selected component is \(\operatorname{Spec}S\).
The selected lift gives \(s_0\), while \(p_1\) gives \(s\).
Both identities in \eqref{eq:ftb-two-point-component-maps}
are built into the construction.
\end{proof}

\begin{srlemma}[Knownness and the actual polarization ambiguity]
\label{lem:ftb-known-same-polarization}
A polarized abelian variety over a number field having the
same principal polarization biextension as a known one is
potentially known.
\end{srlemma}
\begin{proof}
Over the algebraic closure, identify the two abelian varieties.
The quotient of their symmetric principal lines is in
\(\operatorname{Pic}^0[2]\).
The principal polarization identifies it with a \(2\)-torsion
point, so the lines differ by translation by that point.
The point and a line isomorphism descend to one finite
extension, by adjoining their finitely many coordinates.

The line normalization is essential.
For the translated key datum its cube/rigidification data
supplies a unit \(\nu\).
In an old normal form with \(c=1,M=1\), write the point as
\((m+\tau n)/2\).
The half-period theta formula gives the new characteristic
and scalar \(c'\) with
\[
                               (c')^8\nu^2=1.
\]
Explicitly, let \(\lambda\) be the value of the translation
generator at that half period, and let \(d\) be the scalar
in the theta translation formula.
The normalized doubling generator has value \(1\).
Evaluation at the half period and its double gives
\(\lambda^8=\nu^2d^8\), and \(c'=d/\lambda\).
Thus translation multiplies the normal value by \(\nu^2\).
All these identities are over the same finite extension,
and are preserved by every field embedding.

For arbitrary frames the exact change rule is
\[
                    t(a\theta,Bv)=\det(B)^4a^{-8}t(\theta,v).
 \tag{F7}\label{eq:ftb-known-frame-change}
\]
Its multiplier belongs to the field of definition.
Hence every key frame on every further pullback is known.
\end{proof}

\begin{srproposition}[The arithmetic connected-family join]
\label{prop:ftb-connfam-arithmetic-join}
Under the two stated moduli inputs, with their actual period,
level and polarization pullback comparisons, every complex key
datum satisfies \eqref{eq:ftb-connfam}.
\end{srproposition}
\begin{proof}
Fix \(x/\mathbb C\) of dimension \(g\).
The first input gives \(F/K,k,\zeta,\sigma\).
Represent the period matrix and standard level of \(x\)
by \(p:\operatorname{Spec}\mathbb C\to F\).
For the known point use the \(g\)-fold product of
\(E:y^2=x^3-x\), of discriminant \(64\).
The elliptic normalization and product identity give normal
value \(64^g\).
After a finite number-field extension \(K_1/K\), all its
\(k\)-torsion and the compatible level are defined.
Represent it by \(q:\operatorname{Spec}K_1\to F\).
The second input puts its complex pullback and \(p\) on the
same component.

Choose a single affine open \(\operatorname{Spec}R\) containing
both points.  Quasiprojectivity gives it by a hypersurface
avoiding their finite set of images.
It is smooth over \(K\).
The required agreement on idempotents of
\(\mathbb C\otimes_KR\) follows from the ambient component
condition: a smooth Noetherian complex scheme is regular,
so its irreducible components are disjoint; the intersection
of any one with the affine open remains irreducible.
Thus the two points lie on the same component of this
affine open as well.

Apply Lemma~\ref{lem:ftb-symmetric-principal-root-cover} to the
universal biextension on \(R\), obtaining \(R_1,Q\), and lift
\(p\) to \(p_1\).
The fibre of \(Q\) and the line of \(x\) have the same
biextension, so their discrepancy is a \(2\)-torsion translate.
The full \(2\)-torsion is constant over \(R_1\), and the
required point therefore lifts to a section.
Translate \(Q\) by that section and rigidify by its value
there.  The resulting \(Q'\) has the actual line of \(x\)
at \(p_1\), and its principal biextension is unchanged.

Lemma~\ref{lem:ftb-two-point-etale-integral-component} gives
\(k,S,s,s_0\).
The pullback of \(Q'\) has the prescribed polarized fibre
at \(s\); its fibre at \(s_0\) has the biextension of the
known product.
Lemma~\ref{lem:ftb-known-same-polarization} makes that fibre
potentially known.
Make its finite constant-field extension now.
Geometric integrality preserves a smooth integral base;
the complex embedding extends, and both points lift.

The theta sections form an invertible module, by principal
polarization cohomology and base change; the Lie module is
projective of rank \(g\).
At the two primes defined by \(s,s_0\), both are free over
the semilocal ring.
Choose their frames there, clear the finitely many
denominators, and localize by one element outside both
primes.  This gives a key datum \(y\) while retaining both
points.
Before comparing frames, transport the rigidification of \(x\)
through the actual polarized line-bundle pullback isomorphism.
The source's key-on-a-fixed-polarization comparison identifies
that transported rigidification with the chosen one; its line
isomorphism is included in this comparison.  The remaining
parameters are a theta scalar and a Lie matrix.
To impose exactly the specified frames of \(x\), use
\[
                  S'=S[X,X_{ij},(X\det(X_{ij}))^{-1}].
 \tag{F8}\label{eq:ftb-universal-key-frame-torsor}
\]
Multiply the theta section by \(X\), and the Lie frame by
the invertible matrix \((X_{ij})\).
At the complex point take the scalar and matrix carrying
the existing frames to those of \(x\); at the known point
take \(1,I_g\).
The new base is smooth and geometrically integral.
Its complex fibre is the exact key datum \(x\), while the
known fibre is unchanged.
Knownness gives one \(t_0\in k\) at every embedding and for
every corresponding key pullback.
These are precisely the clauses of \eqref{eq:ftb-connfam}.
\end{proof}

\paragraph{Frozen source concordance and the separate lower producers.}
The prefix is
\srcid{src/Development/AbelianVarietyTheory/Proofs/}.
The torsion frame is
\srcid{ThetaLawProofs/ExistsFiniteEtaleCoverConstantTorsion.lean}
(\srcid{f039af8cf45f309722d9c38fdb77e32c50665d19}).
The actual symmetric root uses, in \srcid{ThetaLawProofs},
\srcid{ExistsSymmetricLinearisationSquareDiagonalBundle.lean}
(\srcid{994c2070f67e20bdb9ce54c8efabc9c82511efb9}),
\srcid{ExistsDescentSymmetricLinearisedBundle.lean}
(\srcid{b4d02b09f887aa1bc7e23314b3304db2ee82f9d8}),
\srcid{LamIsoOfDoublingPullbackEqDiagonalSq.lean}
(\srcid{87e1d8620f1761e05f668291e221f8b66f0c781d}),
\srcid{ExistsPrincipalPolarizationSymmetricSquareRoot.lean}
(\srcid{481ca29a56b6d87791856ae5d24b605e3a0cc2d0}),
and \srcid{ExistsSymmetricSqrtPrincipalSiegelFiniteEtale.lean}
(\srcid{ad982d88a3039d9157db361c7eb47c1d2158141f}).
The two-point component argument is
\srcid{GeomIntegralOfEtaleFaithfullyFlatCover.lean}
(\srcid{2c3cdb5d68e977f5fcbac8949e3f558d7ceeb1ce}),
with \srcid{ExistsFiniteExtGeomIntegralLocalization.lean}
(\srcid{ce4a22acec9b2c4be3162269c3732e106c48ecd4}).
Knownness is
\srcid{PotentiallyKnownOfSameLambdaBundle.lean}
(\srcid{46b881f6de2bcd933fdf0af3f175bf70499ce687}),
\srcid{ThetaFrameValueKnownOfTwoTorsionTranslate.lean}
(\srcid{5cbd8eaa06fc6135efd5dcd411fd1f942367eb26}),
\srcid{ThetaNormalFormTranslateHalfPeriod.lean}
(\srcid{72dd6206484af0e01b4c33ad7eecbbb917c36da1}),
\srcid{KeyNormalFormDetRatioTwoTorsionTranslate.lean}
(\srcid{3429a6c1eec5af7b9062c8490ea52e893ec1aa9d}),
and \srcid{ThetaFrameValueKnownOfPolarizedPullback.lean}
(\srcid{391e84ccf8fbd2fe2dfae08549198056345d0a0e}).
The final frame construction is
\srcid{ExistsKeyDataNearTwoPrimes.lean}
(\srcid{b08147dbe77819b1d425e3345501c1fb16b9a1eb})
and \srcid{ExistsThetaKeyFamilyUniversalFrames.lean}
(\srcid{8950e79c3f27fed92a98df44b05825d3e1cfdb39}).
The joins are
\srcid{PolarizedFamilyOfThetaUniformization.lean}
(\srcid{6e9fb2cd94e212f61adaaa66272037b6e4195b01})
and \srcid{KeyDatumConnectedFamilyOfPolarizedArithmeticFami.lean}
(\srcid{3414212fc4b9f57646d6e2b692d061f799a58244}).

The first still separate moduli producer is
\srcid{ModuliProofs/NonemptySiegelFamilyOverArithmeticBase.lean}
(\srcid{c8933879c627986eb6412db2774dfd3f8d998c8b}).
It constructs the fine family through the actual eighth-power
sharp theta closure; generic dimension, special tangent
dimension, universal mapping and principal-Siegel recovery
are substantive lower lemmas.
Its number-field specialization is
\srcid{ThetaLawProofs/ExistsSiegelFamilyOverNumberField.lean}
(\srcid{a87124f2cc26de113f66ea97d5641108f3ff26c9}).
The second input is
\srcid{SiegelFamilyStandardLevelPointsSameComponent.lean}
(\srcid{26a36ef9ff5a0ee06e4f6ba603ea75f78fbaa829}).
Its local producer is
\srcid{SiegelFamilyNearbyPeriodsSameComponent.lean}
(\srcid{0908e63fa80faff82352cd114d3ce644623a521a}),
using actual relative periods, level/Weil compatibility,
uniqueness at a fixed period matrix, and openness of the
locally injective period chart.
These and the fixed-fibre complex realization remain
explicit until their own translations are joined.


\subsection{Specialization and the exponent-one algebraic rule}
\label{subsec:ftb-key-specialization-descent}

In this subsection we assume the exact connected-family and local
analytic family inputs used in
Proposition~\ref{prop:ftb-key-complex-equivariance}, together with
the already constructed normalized powered rule \(\rho\).
Thus \(N(g)\ge1\), \(\rho\) is functorial on characteristic-zero
fields, is a unit on every DVR carrying a key datum, and
\[
                             \rho(x)=t(x)^{N(g)}
                                      \quad(x/\mathbb C).
 \tag{S1}\label{eq:ftb-key-powered-normal-value}
\]
The value \(t(x)=\det(M)^4/c^8\) is independent of normal form and
is equivariant under \emph{all} field maps
\(\mathbb C\to\mathbb C\).
These hypotheses retain the outstanding family producers stated
in the preceding subsection; the argument below does not assert
their existence a second time.

\begin{srlemma}[Countable avoidance and a smooth curve through a point]
\label{lem:ftb-countable-smooth-through}
Let \(R\) be a domain finitely generated over \(\mathbb Z\), with
a map \(a:R\to\mathbb C\).
There are a smooth integral finitely generated
\(\mathbb C\)-algebra \(E\), an injection \(i:R\hookrightarrow E\),
and a complex point \(e:E\to\mathbb C\), with \(ei=a\).
One can take \(\dim E\le1\).
Moreover, if \(E_1\) is any integral finite-type
\(\mathbb C\)-algebra and \(T\subset E_1\setminus\{0\}\) is
countable, some complex point of \(E_1\) is nonzero on all of \(T\).
\end{srlemma}
\begin{proof}
First prove the last assertion.
Noether normalization gives an integral finite inclusion
\(P=\mathbb C[X_1,\ldots,X_d]\hookrightarrow E_1\).
For each nonzero \(r\in T\), integrality and the domain
hypothesis give \(0\ne b_r\in P\cap(r)\).
For example a monic relation for \(r\) of least degree has
nonzero constant term, which belongs to \((r)\).
The nonvanishing set of a nonzero polynomial on \(\mathbb C^d\)
is open and dense: vanishing on an open polydisc would make
the polynomial zero.
The Baire theorem gives \(c\in\mathbb C^d\) at which every
\(b_r\) is nonzero.
A maximal ideal of \(E_1\) above the evaluation ideal at \(c\)
has residue field \(\mathbb C\), by integrality and algebraic
closedness.  Its evaluation cannot kill any \(r\).

To construct the smooth curve, start with
\(\mathbb C\otimes_{\mathbb Z}R\) and the evaluation induced by \(a\).
Flat going down supplies a prime below its evaluation kernel
whose contraction to \(R\) is zero.
Its quotient \(B\) is an integral finite-type
\(\mathbb C\)-algebra, the map \(R\to B\) is injective, and the
given evaluation descends.
It remains to choose a curve through this particular point
without losing that injection.

Choose a finite integral Noether normalization
\(P=\mathbb C[X_1,\ldots,X_d]\hookrightarrow B\).
For every nonzero image of an element of \(R\), choose a
nonzero \(b_r\in P\cap(r)\) as above.
There are only countably many such \(r\).
Let \(c_0\in\mathbb C^d\) be the image of the prescribed point,
and choose \(c_1\) at which all \(b_r\) are nonzero.
Use the line substitution
\[
                        X_j\longmapsto (c_0)_j+
                                      ((c_1)_j-(c_0)_j)T.
 \tag{S2}\label{eq:ftb-line-substitution}
\]
Its kernel is prime and is contained in the contraction of
the prescribed maximal ideal.
Going down for the integral extension from the normal
polynomial ring supplies a prime \(Q\) of \(B\), below that
maximal ideal, lying over this line kernel.
It avoids every \(r\): otherwise \(b_r\) would restrict
identically to zero on the line, contradicting its value
at \(c_1\).
The quotient \(B/Q\) has dimension at most one.
Indeed its contraction quotient is either \(\mathbb C\)
or \(\mathbb C[T]\), and integral incomparability prevents
a longer prime chain upstairs.

Take the integral closure of \(B/Q\) in its fraction field.
It is finite, since \(B/Q\) is a finite-type curve over a field.
It is a Noetherian normal domain of dimension at most one,
hence regular, and is smooth over the perfect field
\(\mathbb C\).
The given point lifts by lying over, with residue field
\(\mathbb C\).
Its coordinate ring is \(E\); the avoidance of all nonzero
\(r\) gives the asserted injection \(R\hookrightarrow E\).
This is the source's line construction and normalization;
it is not an arbitrary smoothing of the original base.
\end{proof}

\begin{srlemma}[Finite-coefficient and countable local-PID descent]
\label{lem:ftb-key-countable-local-pid-specialization}
Let \(V\) be a characteristic-zero local PID domain.
A key datum over \(V\), together with a countable set of its
coefficients, descends to a countable local PID subring
\(V_0\subset V\).
The inclusion is a local homomorphism and is faithfully flat.
If only a finite set of extra coefficients is prescribed and
\(V\) is Noetherian local, the key datum descends to a domain
finitely generated over \(\mathbb Z\), contained in \(V\);
the latter subring need not be local.
\end{srlemma}
\begin{proof}
We first record the ring construction.
For a countable \(T\subset V\), let \(k_0\) be the subfield
of \(\operatorname{Frac}V\) generated by \(T\), and put
\(V_0=V\cap k_0\).
This is countable.
Divisibility of its elements is reflected from \(V\):
if \(b=ac\) in \(V\) and \(a\ne0\), then \(c=b/a\in k_0\),
so \(c\in V_0\).
In particular the inclusion reflects units, and \(V_0\)
is local.
It is a PID.  For an ideal \(I\subset V_0\), choose among
the ideals \(aV\), \(a\in I\), a maximal one, using
Noetherianity of \(V\).
Total divisibility in the local PID \(V\) shows that its
chosen \(a\) divides every \(b\in I\) in \(V\).
Divisibility reflection gives \(I=aV_0\).
The \(V_0\)-module \(V\) is torsion free, hence flat over
this PID, and a flat local inclusion is faithfully flat.

Here is the geometric descent used with this construction.
Choose a very ample projective embedding of the abelian
scheme, from a sufficiently high polarization power.
On the finitely many standard affine charts its defining
ideals have finite sets of polynomial generators.
Adjoin their coefficients.
The closed image in projective space over the resulting
subring pulls back to exactly the original scheme:
on each chart the ideal is the extension of those generators.
Adjoin the finitely many coefficients needed for the
identity, multiplication and inverse maps, and for the
line's finite transition functions and their inverses.
These maps are descended on finite affine refinements;
their polynomial identities are checked after the
faithfully flat inclusion.
Thus the group laws descend.
Flatness, properness, smoothness, the relative dimension,
the line isomorphisms giving symmetry and rigidification,
relative ampleness, and the principal polarization
condition descend by faithful flatness.
This gives the actual polarized abelian scheme over a
countable local PID of the displayed kind.

Choose a key frame on that local base.
Pull it back to \(V\).
The prescribed key datum differs from it by its line
isomorphism, a theta scalar \(a\in V^\times\), and a Lie
matrix \(M\in\operatorname{GL}_g(V)\).
Adjoin the scalar and matrix entries to the countable
coefficient set and repeat the intersection construction.
Unit reflection gives \(a^{-1}\) and
\(\det(M)^{-1}\) in the new local subring.
Change the descended frames by these actual parameters.
The line, its rigidification, the theta section and all
ordered tangent vectors now pull back to the original
key datum, as asserted.

For the finite-type assertion, first descend the polarized
projective data to finitely many coefficients.
Near the contraction of the closed point choose a
nonzero theta section and a Lie basis.
Invert a single element not vanishing there.
Its image is a unit of \(V\); adjoin its inverse, the
required theta scalar and inverse, the Lie matrix entries
and inverse determinant, and the specified finite
coefficient set.
The same actual rescaling gives a key datum over this
finite-type subring.
This is the finite-coefficient construction in the source,
not a claim that an arbitrary frame extends before the
necessary coefficients and inverses are adjoined.
\end{proof}

\begin{srproposition}[Specialization of the unpowered complex value]
\label{prop:ftb-key-normal-value-specializes}
Let \(V\) be a local PID domain with a key datum \(y/V\).
Let \(f,f':V\to\mathbb C\), with \(f\) injective.
If \(u\in V\) and \(t(f^*y)=f(u)\), then
\[
                                t((f')^*y)=f'(u).
 \tag{S3}\label{eq:ftb-key-in-families}
\]
\end{srproposition}
\begin{proof}
Descend \(y,u\) by the finite-type assertion of
Lemma~\ref{lem:ftb-key-countable-local-pid-specialization}, to a finite-type
domain \(R\subset V\) and a key datum \(y_R\).
Apply Lemma~\ref{lem:ftb-countable-smooth-through} to
\(f'|_R\), obtaining \(i:R\hookrightarrow E\) and \(e:E\to\mathbb C\).
Write \(z=i^*y_R\).
The smooth-base powered-unit lemma gives
\(b\in E^\times\) whose evaluation at every complex point is
\(t(z_s)^{N(g)}\).

If \(s:E\to\mathbb C\) does not kill any nonzero \(i(r)\),
then \(si:R\to\mathbb C\) is injective.
The two injections \(f|_R\) and \(si\) admit embeddings
\(\alpha,\beta:\mathbb C\to\mathbb C\) into a common
complex field copy with
\(\alpha f|_R=\beta si\).
This follows by extending the two embeddings of the
countable fraction field, using transcendence bases and
algebraic closures.
Equivariance of \(t\) under these possibly discontinuous
field maps, and its independence of normal form, give
\[
                              t(z_s)=s(i(u)).
\]
Indeed the two key pullbacks become identical after the
common extension, and \(\beta\) is injective.

It follows that \(i(u)^{N(g)}=b\) in \(E\).
If their difference were nonzero, countable avoidance
would produce a complex point not killing that
difference or any nonzero \(i(r)\), contradicting the
last equality and the defining equality for \(b\).
Hence \(i(u)\) is a unit.
Choose one of these generic complex points \(s_0\).
The root-fixed-by-one-point lemma, now on the connected
smooth integral complex curve \(\operatorname{Spec}E\),
compares \(t(z_s)\) with \(s(i(u))\) at every point:
their \(N(g)\)-th powers agree, and their ratio is fixed
by its value \(1\) at \(s_0\).
At \(e\) this gives \(t(z_e)=e(i(u))=f'(u)\).
The fibre \(z_e\) is the prescribed special key pullback,
so this is \eqref{eq:ftb-key-in-families}.
\end{proof}

\begin{srtheorem}[The exponent-one field rule and its DVR units]
\label{thm:ftb-key-rule-exponent-one}
Under the inputs of this subsection there is a unique
rule \(\kappa_K(x)\in K^\times\), for every characteristic-zero
field \(K\) and key datum \(x/K\), which agrees with
\(t(x)\) over \(\mathbb C\), is functorial under all field
maps and all key pullbacks, and has
\[
                   a^8\kappa(a\theta,Mv)=\det(M)^4\kappa(\theta,v).
 \tag{S4}\label{eq:ftb-key-rule-exponent-one-homog}
\]
For every DVR \(V\), fraction field \(K\) of characteristic
zero, and key datum \(y/V\), its generic value is the
image of a unit \(u\in V^\times\).
For every characteristic-zero residue-field quotient
\(r:V\twoheadrightarrow k\),
\[
                                 \kappa_k(r^*y)=r(u).
 \tag{S5}\label{eq:ftb-key-rule-exponent-one-unit}
\]
For a normalized theta function and the ordered tangent
frame \((2\pi i)^{-1}e_j\), the value is \(1\).
\end{srtheorem}
\begin{proof}
We first define the value over a countable field \(k\).
Embed \(k\) in \(\mathbb C\), and pull the key datum to \(x\).
If \(\alpha\) is any automorphism of \(\mathbb C\) fixing \(k\),
apply \eqref{eq:ftb-key-in-families} with the local PID
\(\mathbb C\), \(f=1\), and \(f'=\alpha\).
It gives the transformed value \(\alpha(t(x))\).
Because \(\alpha\) fixes \(k\), cancellation of their common
\(k\)-pullbacks gives an actual key-pullback isomorphism over
\(1_{\mathbb C}\) between \(\alpha^*x\) and \(x\).
Apply the same assertion with \(f'=1\) using that isomorphism;
it gives the value \(t(x)\).
Thus \(\alpha(t(x))=t(x)\).
The fixed field of all such automorphisms is \(k\).
For an element algebraic over \(k\), a distinct conjugate
extends to an automorphism unless the element lies in \(k\),
since the characteristic is zero.
For a transcendental element, changing a transcendence
basis element and extending to algebraic closures gives
a nonfixing automorphism.
Therefore \(t(x)=s(u)\) for a unique \(u\in k\).
Equation \eqref{eq:ftb-key-in-families}, with \(V=k\), makes
this same \(u\) work at every embedding of \(k\).

For an arbitrary \(K\), descend the key datum to a
countable subfield \(k\subset K\), using
Lemma~\ref{lem:ftb-key-countable-local-pid-specialization} and then taking
the fraction field.
Define its value by the preceding \(u\).
This is independent of descent.
Indeed put two countable models in one countable
intermediate field, adjoin the finitely many coefficients
of their common abelian-variety isomorphism and line
isomorphism, and descend these isomorphisms.
Their equality with the given rigidification, section
and tangent vectors descends by faithful flatness.
Embed this common countable field in \(\mathbb C\).
Both proposed values are then the same complex normal
value; injectivity proves equality before embedding.
This also proves functoriality, by pushing one countable
model through a field map.
The complex normal value is nonzero, so the value belongs
to \(K^\times\).
The same descent including \(a,M\) and their inverses
transfers the complex normal-form identity
\eqref{eq:ftb-known-frame-change} to \(K\), giving
\eqref{eq:ftb-key-rule-exponent-one-homog}.

The relation with the powered rule is
\[
                             \kappa_K(x)^{N(g)}=\rho_K(x).
 \tag{S6}\label{eq:ftb-key-root-powered-algebraic}
\]
Check it over a countable model embedded in \(\mathbb C\),
using \eqref{eq:ftb-key-powered-normal-value}, and then
use functoriality.
Over the DVR \(V\), the right side is a unit \(b\in V^\times\).
The element \(\kappa_K(x)\) is integral over \(V\), being
a root of \(T^{N(g)}-b\).
Normality of \(V\) puts it in \(V\); its positive power
being a unit makes it a unit \(u\).

We must also identify the specialization of this chosen
root, not just its power.
Descend \(y,u\) to a countable local PID subring
\(V_0\subset V\), with local faithfully flat inclusion.
The generic countable fraction field embeds in \(\mathbb C\).
By the definition and uniqueness of its value, the
generic complex value is the image of \(u\).
Equation \eqref{eq:ftb-key-in-families} now identifies the
value under every map \(V_0\to\mathbb C\), including maps
factoring through a characteristic-zero residue field.
It therefore represents the value on that residue
pullback in the countable definition.
Pushing to \(k\) proves
\eqref{eq:ftb-key-rule-exponent-one-unit}.
For a normalized theta datum, \(M=I,c=1\), so its value
is \(1\).
All uniqueness assertions follow by the same common
countable model argument.
\end{proof}

\paragraph{Exact scope and source concordance.}
This theorem gives \(\mathrm{Homog}_1\), \(\mathrm{Funct}\),
\(\mathrm{UnitSp}\), and \(\mathrm{ThetaOne}\).
In \(\mathrm{UnitSp}\) the DVR may have positive residue
characteristic; its unit assertion still holds, while the
stated value comparison is only for characteristic-zero
residue fields, exactly as the field rule requires.
The separate archimedean norm identity is not inferred
from the algebraic root argument.

The prefix below is
\srcid{src/Development/AbelianVarietyTheory/Proofs/ThetaLawProofs/}.
The countable avoidance, chosen line and smooth normalization
are \srcid{ExistsAlgHomNeZeroOfCountable.lean}
(\srcid{9dc863ab22573a559ae8c249f33d533a7860cf49}),
\srcid{ExistsDominantComponent.lean}
(\srcid{8e6078832242f5fd87c9969273e88a9b1a4402df}),
\srcid{ExistsPrimeDimensionLEOneAvoiding.lean}
(\srcid{f56c1b74e6815aeea22584eaf146e6252e747ef8}),
and \srcid{ExistsSmoothNormalisation.lean}
(\srcid{2e88aff00b5fb1cfda8566ac7c21761fdce68dff}).
The exact local subring is
\srcid{ExistsCountableLocalPidSubring.lean}
(\srcid{80f4e931782aab0999790894204dc2b851005194}).
Finite projective descent on standard affine charts is
\srcid{VeryAmpleByDescendsCountableSubring.lean}
(\srcid{880d25ecad16a9516141cd62e4c9feee229e7b53}).
Its actual key-data assembly is
\srcid{KeyDataDescendCountableSubrings.lean}
(\srcid{d06ec6cf161d65aff27cc78f9f3693252a1b58d2}),
using finite-coefficient descent for morphisms, line bundles,
group laws and polarizations.
Those are ordinary projective finite-presentation and fpqc
descent operations with the hypotheses stated in the proof;
no uniform bound for their coefficient fields is needed.

The specialization proof, including the finite-type
rescaling of the \emph{actual} key datum, is
\srcid{ExistsSmoothThrough.lean}
(\srcid{fa1c6cdf3595de268b05d93491d66de7c70a4532}).
The common-countable-isomorphism argument, invariant field
definition, root identity and DVR specialization are
\srcid{ExistsCommonCountable.lean}
(\srcid{f748a7e68705d47f6d415d95bf2e20e2370d2577}).
The final wrapper is
\srcid{../UniformizationProofs/KeyFormulaOfConnFam.lean}
(\srcid{c6229bc642171a481eef723c9bdfffc7b9f68a55}).
Its actual connected-family and analytic-continuity
hypotheses remain precisely those retained above.
